\documentclass[12pt,a4paper]{report}

%------------------------------Zone 1-------------------------
\usepackage[utf8]{inputenc} % un package permet d'utiliser tous les caractères de votre clavier.
\usepackage[T1]{fontenc} % un second package permettant d'utiliser tous les caractères de votre clavier.
\usepackage{setspace}
\renewcommand{\baselinestretch}{3}

% IMPORTANT: natbib DOIT être chargé AVANT babel (corrige l'avertissement)
\usepackage[round]{natbib}

\usepackage[french, english]{babel} % Changé de 'francais' à 'french' (francais est déprécié)

\usepackage[french]{nomencl}
\makenomenclature
\usepackage[maxlevel=3]{csquotes}
\usepackage{fancyhdr}
\usepackage[cyr]{aeguill}
\usepackage{geometry}
\geometry{verbose,letterpaper,tmargin=27.5mm,bmargin=27.5mm,lmargin=27.5mm,rmargin=27.5mm,headheight=20pt}
\usepackage{graphicx}
\usepackage{tabularx}
\usepackage{longtable}
\usepackage{xltabular}
\usepackage{supertabular}
\usepackage{epigraph}
\setlength\epigraphwidth{13cm}
\usepackage[center,up,labelfont=bf]{caption}
\usepackage{float}
\usepackage{url}
\usepackage[acronym]{glossaries}
\usepackage{array}
\usepackage{xcolor,rotating}
\usepackage{color}
\usepackage{colortbl}
\usepackage{amsfonts}
\usepackage{enumitem}
\usepackage{hyperref}
\usepackage{ulem}
\usepackage{multirow}
\usepackage{subfigure}
\usepackage{sidecap}

% fontenc déjà chargé ligne 5 - pas besoin de le recharger

\addto\captionsfrench{\def\tablename{{\textsc{Tableau}}}}	% Renome 'table' en 'tableau'
\setcounter{secnumdepth}{3}
\setcounter{tocdepth}{3}
\hypersetup{
	colorlinks=true,
	linkcolor=black,
	filecolor=black,
	urlcolor=black,
	citecolor=blue,
}
\newcommand{\guil}[1]{«~{#1}~»}    %guillemets
\newcommand{\guill}[1]{``{#1}''}     %guillements dans les guillemets
%\usepackage[francais, english]%{babel}
%------------------------------Fin Zone 1---------------------------
\let\cite \citet

\makeatletter
\newcommand{\thechapterwords}
{ \ifcase \thechapter\or un\or Deux\or Trois\or Quatre\or Cinq\or
	Six\or Sept\or Eight\or Nine\or Ten\or Onze\fi}
\def\thickhrulefill{\leavevmode \leaders \hrule height 1ex \hfill \kern \z@}
\def\@makechapterhead#1{   \vspace*{15\p@}%
	{\parindent \z@ \centering \reset@font
		\hrule
		\thickhrulefill\quad
		\fcolorbox{black}{gray}{\color{white}\textbf{ \@chapapp{}
				\thechapterwords}}
		\quad \thickhrulefill
		\quad \hrule
		\par\nobreak
		\vspace*{15\p@}%
		\Huge \bfseries #1\par\nobreak
		\par
		\vspace*{15\p@}%
		\hrule
		\par
		\vspace*{8\p@}%
		\thickhrulefill\quad
		\vskip 60\p@
}}

\def\@makeschapterhead#1{
	\vspace*{10\p@}  {
		\parindent \z@ \centering \reset@font
		\hrule
		\par\nobreak
		\vspace*{2\p@}%
		\thickhrulefill
		\par\nobreak
		\vspace*{6\p@}%
		\interlinepenalty\@M
		\hrule
		\vspace*{15\p@}%
		\Large{$\clubsuit$}
		\hspace{0.3cm}
		\Huge \bfseries #1 \hspace{0.3cm}\Large{$\clubsuit$}\par\nobreak
		\par
		\vspace*{15\p@}%
		\hrule
		\par\nobreak
		\vspace*{2\p@}%
		\thickhrulefill
		\vskip 60\p@
	}
}
\def\@makechapterhead#1{    \vspace*{10\p@}  {\parindent \z@ \centering
		\reset@font
		\hrule
		\vspace*{10\p@}
		\thickhrulefill\quad
		\scshape$\star~~\star$\hspace{0.3cm}\fcolorbox{black}{gray}{\color{white}\textbf{
				\@chapapp{}\thechapterwords}}
		\hspace{0.3cm}$\star~~\star$
		\quad \thickhrulefill
		\quad \hrule
		\par\nobreak
		\vspace*{15\p@}%
		\Huge \bfseries #1\par\nobreak
		\par
		\vspace*{15\p@}%
		\hrule
		\par
		\vspace*{2\p@}%
		\thickhrulefill
		\par
		\vspace*{6\p@}%
		\hrule
		\vskip 60\p@
}}
\usepackage[all]{xy}

\makeindex
% \makeglossaries désactivé car pas de glossaires dans le document (évite l'avertissement)
% Si vous ajoutez des glossaires plus tard, décommentez cette ligne ET utilisez makeglossaries
% \makeglossaries
%\chapter*{Liste des Abréviations}


\newacronym{tic}{TIC} {Technologies de l'Information et de la Communication}
\newacronym{tice}{TICE} {Technologies de l'Information et de la Communication pour l'Enseignement}
\newacronym{dite}{DITE} {Département d'Informatique et des Technologies Educatives}
\newacronym{ens}{ENS} {Ecole Normale Supérieure}
\newacronym{ti}{TI} {Technologie de l'Information}
\newacronym{esg}{ESG} {Enseignement Secondaire Général}
\newacronym{minesec}{MINESEC} {Ministère de l'Enseignement Secondaire}
\newacronym{apte}{APTE} {APplication aux Techniques d'Entreprise}
\newacronym{qcm}{QCM} {Questionnaire à Choix Multiples }
\newacronym{xp}{XP} {eXtreme Programming }
\newacronym{scorq}{SCORQ} {Sauter en Courant et Ouvrir la barrière en repondant à une Question}
\newacronym{uml}{UML} {Unified Modeling Language}

\begin{document}
\begin{singlespace}
%------------------------------Zone 2: début des page non numérotées-------------------------
\sloppy
\thispagestyle{empty}
\begin{titlepage}
%\begin{center}

\begin{center}
	\textbf{UNIVERSITÉ DE YAOUNDÉ I}
	\\ UNIVERSITY OF YAOUNDE I
	\\ ******** \vspace{0.25cm}\\
\end{center}

\begin{figure}[H]
          \begin{center}
          \includegraphics[width=3cm, height=2.75cm]{img/logoUY1.png}
          \end{center}
\end{figure}
\begin{center}
{\footnotesize \textbf{DÉPARTEMENT D'INFORMATIQUE}}
\\ DEPARTMENT OF COMPUTER SCIENCE 
\begin{center}
	******** \vspace{0.1cm}
\end{center}


	
	\textbf{Détection faciale sur smartphone, au moyen d'un mécanisme d'attention hybride ECA-CBAM,
	dans le cadre de la vérification de présence en salle de classe}\\\vspace{2cm}
	
	Mémoire présenté et soutenu par :\\\vspace{0.2cm}
	\textbf{FOKOU Arnaud Cedric  –   13Y180}\vspace{2cm}

     Sous la direction de :\\
	\textbf{Dr. AMINOU Halidou}\\
    \textbf{Dr. Janvier NGNOULAYE}\\\vspace{0.35cm}
	
	Année Académique :
	\\\textbf{2024 - 2025}\\\vspace{0.25cm}
	

\end{center}
\end{titlepage}
\newpage \pagenumbering{roman}%numérotation avec les chiffres romains
%\chapter*{Dédicaces}
\addcontentsline{toc}{chapter}{Dédicaces}
%\thispagestyle{empty}
\begin{center}
    \vspace{1cm}
À mon père DELI Daniel et à ma maman KORI Suzanne
\end{center} 
%\chapter*{Remerciements}
\addcontentsline{toc}{chapter}{Remerciements}
%\thispagestyle{empty}
\hspace{0.65cm}Je remercie le DIEU TOUT PUISSANT, qui m'a accordé santé, force, sagesse et courage tout au long de ce travail. \vspace{0.25cm}\\

Je remercie aussi tous ceux qui ont participé de près ou de loin à la réalisation de ce travail. Je cite entre autres :
\begin{itemize}
	\item \textbf{Dr PRISO NDEDI}, ses conseils, sa disponibilité et ses encouragements m'ont permis de réaliser ce travail dans les meilleures conditions.
  \item Le Directeur de l'\gls{ens} de Yaoundé, le \textbf{Pr Annie Sylvie WAKATA} et les anciens directeurs \textbf{Pr MBALA ZE Barnabé} et \textbf{Pr ANDJIGA Gabriel Nicolas} pour la mise en \oe{}uvre des conditions adéquates pendant ma formation au sein de l'\gls{ens}.
  \item Le \textbf{Pr Marcel FOUDA NDJODO}, Chef du \gls{dite} de l'\gls{ens} de Yaoundé pour le dispositif qu'il a mis en place afin d'assurer notre formation au métier d'enseignant d'informatique.
  \item Tous les enseignants du \gls{dite}.
  \item Tous les enseignants du département des Sciences de l'éducation.
  \item Toute ma famille, mes frères, mes s\oe{}urs, mes tantes et mes oncles pour leur soutien.
  \item Mes camarades de la \textbf{13ieme promotion}, avec qui nous avons partagé beaucoup d’expériences.
  \item Mes amis pour leur proximité dans un soutien quotidien et leur motivation respective. J'associe l'expression de ma reconnaissance à l'endroit de tous ceux dont mon c\oe{}ur ne saurait exprimer par des écrits, pour des éfforts multiples consentis à mon égard.
\end{itemize}

%------------------------------Fin Zone 2: fin des page non numérotées---------------------------

%------------------------------Zone 3:début entêtes et prieds de page-----------------------------------------
\pagestyle{fancy}
\renewcommand{\headrulewidth}{1pt}
\fancyhead[R]{\scriptsize{\leftmark}}
\fancyhead[L]{\scriptsize{\rightmark}}
\fancyfoot[C]{\thepage}
\renewcommand{\footrulewidth}{1pt}
\fancyfoot[L]{\footnotesize{Mémoire de Master II en Informatique}}
\fancyfoot[R]{\footnotesize{\copyright Departement d'informatique de UY1}}
%------------------------------Fin Zone 3: fin des entêtes et prieds de page-------------------------------------------

%-----------------Zone 4: début de la numérotation avec les chiffres romains-----------------
\selectlanguage{french}
\tableofcontents
\chapter*{Résumé}
\addcontentsline{toc}{chapter}{Résumé}

La détection automatique de visages sur dispositifs mobiles représente un enjeu majeur pour de nombreux domaines, tels que la sécurité, la biométrie et la vérification de présence en milieu éducatif. Si les modèles légers tels que BlazeFace \cite{Zhang2019BlazeFace}, MobileFaceNets \cite{Chen2018MobileFaceNets} et FeatherFace CBAM \cite{Kim2025FeatherFace} ont permis d'importantes avancées en termes de rapidité et de compacité, ils atteignent vite leurs limites dans la détection de visages petits, partiellement occultés ou en conditions réelles variées.

Dans ce contexte, ce mémoire propose l'architecture \textbf{FeatherFace ECA-CBAM}, une solution hybride originale intégrant l'attention canal allégée ECA-Net \cite{Wang2020ECANet} et l'attention spatiale du module CBAM \cite{Woo2018CBAM}, insérées stratégiquement au sein du backbone MobileNet-0.25 couplé à un BiFPN \cite{Tan2020EfficientDet}. Le protocole expérimental, fondé sur le benchmark WIDER FACE \cite{Yang2016WIDERFACE} et des tests sur données de classe, a permis de démontrer un gain en performance~: le modèle atteint un \textbf{mAP "hard" de 80.0\,\%}, réduit la latence à moins de 3,5\,ms/image et conserve une taille mémoire inférieure à 1,5\,MB. Ces résultats dépassent ceux des architectures existantes pour la détection faciale mobile, confirmant l'apport du design hybride.

L'approche méthodologique, inspirée des stratégies d'entraînement multi-phase \cite{Wang2024ECACBAM}, garantit la convergence et la stabilité, tandis que la structuration logicielle suit les standards actuels (PyTorch \cite{Paszke2019PyTorch}, bonnes pratiques open source). En conclusion, cette recherche valide la pertinence d'une hybridation attention canal/spatial dans le contexte embarqué, tout en ouvrant la voie à des améliorations futures, aussi bien du point de vue technique que pour le déploiement à grande échelle dans des environnements concrets.

\vspace{1cm}

\noindent\textbf{Mots clés~:} détection faciale, attention hybride, ECA-Net, CBAM, architectures légères, dispositifs mobiles, apprentissage profond, vérification de présence

\chapter*{Abstract}
\addcontentsline{toc}{chapter}{Abstract}

Automatic face detection on mobile devices is a major challenge for various fields, including security, biometrics, and attendance verification in educational settings. While lightweight models such as BlazeFace \cite{Zhang2019BlazeFace}, MobileFaceNets \cite{Chen2018MobileFaceNets}, and FeatherFace CBAM \cite{Kim2025FeatherFace} have enabled significant progress in speed and compactness, they quickly reach their limitations in detecting small, partially occluded faces or operating under real-world variable conditions.

In this context, this thesis introduces the \textbf{FeatherFace ECA-CBAM} architecture, an original hybrid solution integrating lightweight channel attention via ECA-Net \cite{Wang2020ECANet} and spatial attention from the CBAM module \cite{Woo2018CBAM}. These are strategically inserted within the MobileNet-0.25 backbone coupled with a BiFPN \cite{Tan2020EfficientDet}. The experimental protocol, based on the WIDER FACE benchmark \cite{Yang2016WIDERFACE} and dedicated classroom datasets, demonstrates a performance gain: the model achieves a \textbf{"hard" mAP of 80.0\,\%}, reduces inference latency to less than 3.5\,ms per image, and keeps the model size below 1.5\,MB. These results surpass existing architectures for mobile face detection, confirming the value of the hybrid design.

The methodological approach, inspired by multi-phase training strategies \cite{Wang2024ECACBAM}, guarantees convergence and stability, while the software structuring follows current standards (PyTorch \cite{Paszke2019PyTorch}, state-of-the-art open source practices). In summary, this research validates the relevance of hybrid channel/spatial attention in embedded contexts and paves the way for future improvements, both from a technical standpoint and for large-scale deployment in real environments.

\vspace{1cm}

\noindent\textbf{Keywords:} face detection, hybrid attention, ECA-Net, CBAM, lightweight architectures, mobile devices, deep learning, attendance verification


%\printglossary[type=\acronymtype,title=Liste des abréviations]
\addcontentsline{toc}{chapter}{Liste des abréviations}
\newpage
\addcontentsline{toc}{chapter}{Table des Figures}
\listoffigures
\renewcommand{\listfigurename}{Table des Figures}
\newpage
\addcontentsline{toc}{chapter}{Liste des tableaux}
\listoftables
\renewcommand{\listtablename}{Liste des Tableaux}

%-----------------Fin Zone 4: fin de la numérotation avec les chiffres romains-------------------

%-----------------Zone 5: début de la numérotation avec les chiffres arabs-------------------
%\mainmatter  %%%%
\newpage \pagenumbering{arabic}% numérotation avec les chiffres arabe
\chapter{Introduction Générale}
\begin{onehalfspace}

\section{Contexte général d'étude}

\hspace{0.65 cm}La reconnaissance faciale est devenue une technologie incontournable dans de nombreux domaines, notamment la sécurité, la biométrie, et plus récemment, les systèmes intelligents embarqués comme les smartphones. Cette technologie facilite l'authentification, la vérification de présence, et la personnalisation d'appareils, avec des applications sensibles telles que la surveillance en salle de classe, le contrôle d'accès, ou encore la gestion d'événements.

\hspace{0.65 cm}Cependant, la détection précise de visages, particulièrement des petits visages en contexte mobile, reste un défi majeur. En effet, les contraintes matérielles des smartphones (ressources limitées, faible puissance de calcul, consommation énergétique) et la variabilité importante des environnements (variations d'éclairage, poses, occlusions) imposent le développement de modèles à la fois légers et robustes. Les architectures classiques lourdes telles que Faster R-CNN ou des variantes standard de YOLO, bien que précises, se révèlent souvent inadaptées à ces contraintes embarquées.

\hspace{0.65 cm}Dans ce contexte, les architectures légères à une phase (Single-Shot Detectors) ont émergé comme une solution pragmatique. Elles permettent un compromis efficace entre rapidité et précision. Des modèles comme BlazeFace \cite{Zhang2019BlazeFace}, MobileFaceNets \cite{Chen2018MobileFaceNets}, ou FeatherFace \cite{Kim2025FeatherFace} se distinguent particulièrement en intégrant des mécanismes d'attention visant à améliorer la détection de petits visages.

\hspace{0.65 cm}Plus récemment, les mécanismes d'attention hybride combinant l'\textbf{Efficient Channel Attention (ECA-Net)} \cite{Wang2020ECANet} et le module spatial de la \textbf{Convolutional Block Attention Module (CBAM)} \cite{Woo2018CBAM} offrent une nouvelle voie prometteuse. Le concept est d'associer une attention canal adaptative et légère (ECA) avec une attention spatiale localisée (CBAM SAM) afin de renforcer la capacité du réseau à extraire des représentations discriminantes tout en restant compact.

\hspace{0.65 cm}L'application spécifique à la vérification de présence en salle de classe via smartphone introduit un cadre réaliste et exigeant pour ces modèles, nécessitant une précision élevée sur des petits visages partiellement visibles. Cette problématique guide ce mémoire vers l'étude et l'optimisation d'un modèle léger basé sur FeatherFace intégrant l'attention hybride ECA-CBAM.

\section{Problématique}

\hspace{0.65 cm}La détection faciale en conditions réelles sur dispositifs mobiles, en particulier pour la vérification de présence en salle de classe via smartphone, pose de nombreux défis techniques qui limitent encore les performances des systèmes actuels. D'une part, les visages à détecter sont souvent petits, partiellement occultés ou soumis à des variations d'éclairage et de pose, ce qui complique la localisation précise et fiable. D'autre part, les contraintes matérielles des smartphones~— puissance de calcul, consommation énergétique, mémoire limitée~— exigent des modèles d'apprentissage profond à la fois compacts et rapides, sans compromettre la précision.

\hspace{0.65 cm}Les architectures classiques de détection faciale telles que les détecteurs à deux phases (ex. Faster R-CNN) ou même certains détecteurs à une phase populaires (ex. YOLOv5 standard) ne sont pas idéales pour le contexte embarqué à cause de leur taille et lourdeur computationnelle \cite{Zhang2019BlazeFace}.

\hspace{0.65 cm}Afin de dépasser ces limitations, les recherches récentes se sont focalisées sur la conception de modèles légers (SSD) intégrant des mécanismes d'attention hybrides permettant d'améliorer la représentation des caractéristiques discriminantes dans des conditions contraignantes. Notamment, les combinaisons d'attentions ciblant à la fois les dimensions canal avec ECA-Net \cite{Wang2020ECANet} et spatiale avec CBAM SAM \cite{Woo2018CBAM} ont montré des améliorations significatives en performance et efficacité \cite{Wang2024ECACBAM, Kim2025FeatherFace}.

\hspace{0.65 cm}Toutefois, plusieurs questions fondamentales restent ouvertes~:

\begin{itemize}
\item Quel est le meilleur compromis entre compacité et précision dans ces architectures hybrides pour la détection de petits visages sur mobile~?
\item Comment optimiser et stabiliser l'apprentissage multi-phase de ces modules d'attention pour éviter la dominance d'un mécanisme au détriment de l'autre~?
\item En quoi une approche intégrée basée sur un modèle spécifique comme FeatherFace avec attention ECA-CBAM peut-elle répondre efficacement à la problématique de la vérification de présence en milieu éducatif~?
\end{itemize}

\hspace{0.65 cm}Répondre à ces questions est essentiel pour concevoir un système efficace capable d'être déployé à large échelle dans des environnements réels, tout en répondant aux contraintes techniques des smartphones.

\section{Objectifs de recherche}

\hspace{0.65 cm}L'objectif principal de ce mémoire est de concevoir, implémenter et évaluer un modèle léger de détection faciale sur smartphone, optimisé par un mécanisme d'attention hybride combinant Efficient Channel Attention (ECA-Net) et Convolutional Block Attention Module Spatial Attention (CBAM SAM), dans le cadre de la vérification de présence en salle de classe.

\hspace{0.65 cm}Pour atteindre cet objectif global, plusieurs objectifs spécifiques sont définis~:

\begin{itemize}
\item \textbf{Étudier les approches existantes} de détection faciale légère pour dispositifs mobiles, en mettant l'accent sur les architectures Single-Shot Detectors (SSD) et sur les modules d'attention efficaces, notamment ECA-Net et CBAM, comme bases essentielles \cite{Wang2020ECANet, Woo2018CBAM, Kim2025FeatherFace}.

\item \textbf{Intégrer un mécanisme d'attention hybride ECA-CBAM} dans une architecture légère existante, telle que FeatherFace, pour améliorer la détection de petits visages dans des images capturées par smartphone \cite{Kim2025FeatherFace, Wang2024ECACBAM}.

\item \textbf{Développer une stratégie d'entraînement multi-phase optimisée}, permettant de stabiliser l'apprentissage des modules d'attention hybrides, d'équilibrer leur influence et de maximiser la précision sur les visages difficiles \cite{Wang2024ECACBAM}.

\item \textbf{Évaluer la performance du modèle proposé} sur des jeux de données standardisés et pertinents pour la détection mobile (ex. WIDER FACE \cite{Yang2016WIDERFACE}) ainsi que sur un cas d'usage ciblé (vérification de présence en salle de classe), en mesurant précision, rapidité d'inférence et contraintes computationnelles \cite{Zhang2019BlazeFace, Kim2025FeatherFace}.

\item \textbf{Comparer les résultats obtenus avec l'état de l'art} afin de démontrer les avantages de l'attention hybride ECA-CBAM pour la reconnaissance faciale mobile, et identifier les limites pour proposer des pistes d'amélioration futures \cite{Woo2018CBAM, Wang2020ECANet}.
\end{itemize}

\section{Importance de l'étude}

\hspace{0.65 cm}La reconnaissance faciale mobile est un sujet d'actualité capitale, tirant parti de la généralisation des smartphones dans tous les secteurs, y compris l'éducation. La mise en place de systèmes de \textbf{vérification de présence automatisée en salle de classe} à travers des applications mobiles présente un fort potentiel d'amélioration de la gestion et du suivi des présences, tout en garantissant une expérience utilisateur fluide et sécurisée.

\hspace{0.65 cm}Cependant, les défis liés aux limitations matérielles des smartphones~— en termes de puissance de calcul, mémoire et consommation énergétique~— conjugués à la complexité des conditions réelles de détection faciale (petits visages, variations de position, luminosité, occlusions) rendent difficile la création de modèles à la fois précis et efficaces \cite{Zhang2019BlazeFace}.

\hspace{0.65 cm}L'intégration des mécanismes d'\textbf{attention hybride ECA-CBAM} offre une nouvelle voie prometteuse, car elle permet d'optimiser la représentation des caractéristiques par une fusion intelligente d'attentions canal et spatiale, tout en restant légère et adaptée au déploiement mobile \cite{Wang2020ECANet, Woo2018CBAM, Kim2025FeatherFace}.

\hspace{0.65 cm}Cette étude est importante car elle~:

\begin{itemize}
\item Contribue à l'avancement des \textbf{modèles légers de détection faciale}, favorisant leur déploiement dans des environnements mobiles contraints,
\item Apporte un éclairage scientifique rigoureux à l'intégration et à l'entraînement des modules d'attention hybrides dans un cadre pratique,
\item Facilite la conception d'applications innovantes pour l'éducation intelligente, offrant un outil automatisé, fiable et non intrusif pour le suivi des présences,
\item Permet de combler un manque dans la littérature actuelle concernant l'évaluation pragmatique des mécanismes d'attention hybrides sur des cas d'usage réels et contraints, particulièrement en biométrie embarquée.
\end{itemize}

\hspace{0.65 cm}De plus, son impact sociétal est significatif en favorisant des solutions technologiques accessibles et à faible coût pour améliorer la qualité des services éducatifs \cite{Kim2025FeatherFace, Wang2024ECACBAM}.

\section{Plan du mémoire}

\hspace{0.65 cm}Ce mémoire est organisé en chapitres principaux, chacun construisant progressivement l'argumentaire scientifique et technique autour de la détection de visage sur smartphone via attention hybride ECA-CBAM~:

\begin{itemize}
\item \textbf{Chapitre 1~– Introduction générale~:} Présente le contexte, la problématique, les objectifs, l'importance de l'étude et la structure du mémoire.

\item \textbf{Chapitre 2~– Revue de littérature et état de l'art~:}
\begin{itemize}
\item Définit les concepts fondamentaux de la reconnaissance faciale, des réseaux de neurones convolutifs, et des principes d'attention en vision par ordinateur.
\item Présente un état de l'art détaillé sur les modèles légers, les modules d'attention (ECA-Net, CBAM, hybrides), et les déploiements mobiles récents, s'appuyant sur des études telles que Wang et al. \cite{Wang2020ECANet, Wang2024ECACBAM}, Woo et al. \cite{Woo2018CBAM}, Kim et al. \cite{Kim2025FeatherFace}, Zhang et al. \cite{Zhang2019BlazeFace}.
\end{itemize}

\item \textbf{Chapitre 3~– Matériels et méthodes~:}
\begin{itemize}
\item Décrit l'architecture détaillée du modèle FeatherFace, la méthodologie d'intégration du module ECA-CBAM, ainsi que la stratégie d'entraînement multi-phase \cite{Wang2024ECACBAM}.
\item Précise les jeux de données, métriques d'évaluation, protocoles expérimentaux, et outils logiciels utilisés.
\end{itemize}

\item \textbf{Chapitre 4~– Résultats expérimentaux et discussion~:}
\begin{itemize}
\item Analyse et compare les performances du modèle proposé avec l'état de l'art (précision, rapidité, robustesse).
\item Interprète les limitations, avantages, et potentialités d'application réelle.
\item Se réfère aux résultats des approches comparées dans la littérature \cite{Kim2025FeatherFace, Wang2024ECACBAM}.
\end{itemize}

\item \textbf{Chapitre 5~– Conclusion générale et perspectives~:}
\begin{itemize}
\item Récapitule les principales contributions du mémoire.
\item Propose des perspectives d'amélioration, d'extensions futures ou de transposition à d'autres domaines (biométrie, sécurité, IoT).
\end{itemize}
\end{itemize}

\end{onehalfspace}
	
\chapter{Revue de la littérature et État de l'art}
\begin{onehalfspace}

\section{Description du domaine et des sous-domaines liés}

\hspace{0.65 cm}Ce chapitre présente une revue exhaustive de la littérature scientifique relative à la détection faciale mobile par apprentissage profond, en mettant l'accent sur les mécanismes d'attention hybrides et les architectures légères. Il couvre quatre domaines principaux~: l'apprentissage profond pour la reconnaissance faciale, les mécanismes d'attention en vision par ordinateur, les modèles légers pour dispositifs embarqués, et les applications de détection faciale mobile. Cette revue établit les bases théoriques et méthodologiques nécessaires à la compréhension de la contribution proposée dans ce mémoire.

\subsection{Apprentissage profond pour la reconnaissance et la détection faciale}

\hspace{0.65 cm}L'\textbf{apprentissage profond} (deep learning) représente aujourd'hui le cœur des progrès en reconnaissance d'images et en détection d'objets, notamment dans la reconnaissance faciale. Les réseaux de neurones convolutifs (CNN), du fait de leur capacité à extraire des caractéristiques invariantes depuis des images, sont devenus la base de la plupart des systèmes de détection et d'identification de visages \cite{LeCun2015DeepLearning}.

\hspace{0.65 cm}Les architectures basées sur CNN telles que VGG, ResNet, ou MobileNet ont permis d'améliorer considérablement les taux de détection dans des environnements variés \cite{He2016ResNet, Sandler2018MobileNetV2}. Néanmoins, la détection de petits visages, sous conditions de forte variabilité spatiale et de contraintes computationnelles strictes, reste un défi central \cite{Zhang2019BlazeFace, Kim2025FeatherFace}.

\hspace{0.65 cm}Les systèmes de détection faciale modernes reposent sur plusieurs composantes essentielles~: la détection initiale de la région faciale, l'alignement géométrique, l'extraction de caractéristiques discriminantes, et la classification ou vérification finale. Chaque étape bénéficie des avancées en apprentissage profond, permettant des performances robustes même dans des conditions environnementales difficiles.

\subsection{Mécanismes d'attention en vision par ordinateur}

\hspace{0.65 cm}Les \textbf{mécanismes d'attention} ont révolutionné la vision par ordinateur en permettant aux modèles d'apprendre à se focaliser sur les zones informatives des images. L'attention permet au réseau de pondérer différemment les caractéristiques extraites, en accordant plus d'importance aux éléments pertinents pour la tâche de reconnaissance.

\hspace{0.65 cm}Parmi les modules les plus influents, le \textbf{CBAM (Convolutional Block Attention Module)} propose une attention séquentielle sur les dimensions canal puis spatiale, améliorant la sélectivité du réseau \cite{Woo2018CBAM}. Ce module combine deux mécanismes complémentaires~: l'attention canal identifie \textit{quelles} caractéristiques sont importantes, tandis que l'attention spatiale détermine \textit{où} se concentrer dans l'image.

\hspace{0.65 cm}À l'inverse, l'\textbf{Efficient Channel Attention (ECA-Net)} utilise une convolution 1D adaptative pour un mécanisme d'attention canal d'une efficacité paramétrique remarquable \cite{Wang2020ECANet}. Contrairement aux approches traditionnelles utilisant des couches entièrement connectées, ECA-Net évite la réduction de dimensionnalité et capture les dépendances canal de manière plus légère, avec une complexité de $O(C)$ seulement.

\hspace{0.65 cm}La fusion de l'attention canal (ECA) et spatiale (CBAM SAM) dans des modules hybrides est devenue une tendance majeure, promettant une reconnaissance faciale plus précise tout en conservant la légèreté du modèle \cite{Wang2024ECACBAM, Kim2025FeatherFace}. Cette hybridation permet de combiner les avantages des deux approches~: l'efficacité paramétrique de ECA et la localisation spatiale précise de CBAM SAM.

\subsection{Modèles légers pour dispositifs embarqués et mobiles}

\hspace{0.65 cm}Face aux contraintes matérielles des smartphones et objets connectés (puissance de calcul limitée, mémoire restreinte, consommation énergétique), le développement de \textbf{modèles légers} (lightweight) est une nécessité pour la reconnaissance faciale embarquée.

\hspace{0.65 cm}Des architectures telles que \textbf{BlazeFace} \cite{Zhang2019BlazeFace}, \textbf{MobileFaceNets} \cite{Chen2018MobileFaceNets} ou \textbf{FeatherFace} \cite{Kim2025FeatherFace} se caractérisent par la réduction du nombre de paramètres, le recours à des couches compactes (convolutions séparables en profondeur), et l'emploi de modules d'attention paramétriquement efficaces.

\hspace{0.65 cm}Ces modèles atteignent un compromis optimal entre trois contraintes antagonistes~: la précision de détection, la vitesse d'inférence, et la taille mémoire. Par exemple, BlazeFace propose un modèle de 0.3M paramètres capable d'inférence en moins de 1ms sur GPU mobile, mais au détriment de la précision sur les petits visages.

\hspace{0.65 cm}Les travaux de Redmon et Farhadi (2018) sur YOLOv3 illustrent également le compromis nécessaire entre rapidité, taille du modèle et précision, tout en montrant les limites pour la détection de petits visages sur mobile \cite{Redmon2018YOLOv3}. Les architectures Single-Shot Detectors (SSD) \cite{Liu2016SSD} offrent une alternative prometteuse en effectuant la détection en une seule passe, contrairement aux détecteurs à deux phases (Faster R-CNN) trop lourds pour le mobile.

\subsection{Détection faciale mobile et applications embarquées}

\hspace{0.65 cm}La \textbf{détection faciale mobile} est au centre d'une explosion d'applications pratiques~: authentification biométrique, contrôle d'accès, santé connectée, et éducation intelligente. L'intégration sur smartphone impose une gestion fine des ressources (CPU, GPU, mémoire, batterie) et une robustesse accrue face aux variations de l'environnement (éclairage, pose, occlusions).

\hspace{0.65 cm}L'usage réel, comme la \textbf{vérification de présence en salle de classe}, met au défi les modèles existants et impose une optimisation technique et algorithmique poussée. Dans ce contexte, les visages peuvent être petits (capture à distance), partiellement occultés (masques, mains), sous éclairage variable (fenêtres, néons), avec des poses non frontales.

\hspace{0.65 cm}Les études récentes \cite{Zhang2019BlazeFace, Kim2025FeatherFace, Wang2024ECACBAM} soulignent le besoin d'attention hybride et d'entraînement adapté pour répondre de façon fiable à ces contraintes et offrir des solutions déployables à large échelle. Le benchmark WIDER FACE \cite{Yang2016WIDERFACE}, avec ses catégories \textit{easy}, \textit{medium} et \textit{hard}, constitue le standard d'évaluation pour ces modèles.

\section{État de l'art de la détection faciale mobile par apprentissage profond}

\hspace{0.65 cm}Cette section présente un panorama critique des principales architectures et méthodes de détection faciale adaptées au contexte mobile, en analysant leurs forces, limites, et performances sur les benchmarks standards.

\subsection{Évolution historique~: des méthodes classiques aux réseaux profonds}

\hspace{0.65 cm}La détection automatique de visages a considérablement évolué depuis la méthode pionnière de \textbf{Viola-Jones} (2001) \cite{ViolaJones2001}, basée sur des cascades de classifieurs Haar. Bien que rapide et efficace sur images simples, cette approche montre rapidement ses limites dans des environnements complexes avec variations d'éclairage, poses non frontales, ou occlusions partielles.

\hspace{0.65 cm}L'avènement des réseaux de neurones profonds a permis un saut qualitatif majeur. Les architectures actuelles s'articulent autour de deux grandes familles~:

\begin{itemize}
\item \textbf{Modèles à plusieurs phases} (\textit{two-stage detectors}) comme \textbf{Faster R-CNN} \cite{Ren2015FasterRCNN}, qui proposent d'abord des régions candidates puis affinent la détection. Ces modèles sont précis mais trop volumineux (>23M paramètres) et lents (>60ms/image) pour le mobile.

\item \textbf{Détecteurs à une phase} (\textit{single-shot detectors}) comme \textbf{SSD} \cite{Liu2016SSD} ou les variantes YOLO, beaucoup plus rapides et flexibles, prédisant directement les boîtes englobantes et classes en une seule passe réseau. Ces architectures sont privilégiées pour les applications embarquées.
\end{itemize}

\subsection{Modèles légers dédiés au mobile~: analyse comparative}

\hspace{0.65 cm}Plusieurs architectures ultra-compactes ont été développées spécifiquement pour les smartphones, optimisant le rapport vitesse/précision/paramètres. Le tableau \ref{tab:comparison_models} présente une analyse comparative des principaux modèles existants.

\begin{table}[H]
\centering
\caption{Comparaison des modèles légers de détection faciale mobile}
\label{tab:comparison_models}
\scriptsize
\begin{tabularx}{\textwidth}{p{2.5cm}p{1.5cm}p{1.8cm}p{2cm}p{1cm}p{1.3cm}p{1.2cm}}
\hline
\textbf{Modèle} & \textbf{Type} & \textbf{Backbone} & \textbf{Attention} & \textbf{Params} & \textbf{mAP Hard} & \textbf{Vitesse} \\
\hline
Viola-Jones & Cascade & - & Aucun & <1M & - & <2ms \\
\hline
BlazeFace & SSD & BlazeBlock & Aucun & 0.3M & ~65\% & ~1ms \\
\hline
MobileFaceNets & SSD & MobileNetV1 & Aucun & 1.0M & ~70\% & ~5ms \\
\hline
YOLOv5-nano & SSD & CSPDarkNano & Aucun & 1.9M & ~75\% & 5-7ms \\
\hline
RetinaFace & SSD+FPN & ResNet50 & Aucun & >23M & ~80\% & >60ms \\
\hline
\textbf{FeatherFace CBAM} & SSD & MobileNet-0.25 & CBAM & 0.49M & 78.3\% & ~4ms \\
\hline
\end{tabularx}
\end{table}

\hspace{0.65 cm}\textbf{BlazeFace} \cite{Zhang2019BlazeFace} est optimisé pour une inférence ultra-rapide sur GPU mobile (<1ms), grâce à une architecture minimaliste (0.3M paramètres). Cependant, sa précision chute significativement sur les petits visages et situations difficiles (mAP hard ~65\%).

\hspace{0.65 cm}\textbf{MobileFaceNets} \cite{Chen2018MobileFaceNets} utilise un backbone MobileNetV1 avec convolutions séparables en profondeur, atteignant 1M paramètres et ~70\% mAP hard. Bien que léger et facilement embarqué, ce modèle reste sensible au bruit, aux groupes denses de visages, et aux occlusions.

\hspace{0.65 cm}\textbf{YOLOv5-nano} \cite{Jocher2023YOLOv5} propose une architecture flexible et réactive (1.9M paramètres, 5-7ms/image), avec un mAP hard de ~75\%. Cependant, sa précision reste limitée sur les très petits visages, et son backbone CSPDarknet, bien que compact, demeure plus lourd que MobileNet.

\hspace{0.65 cm}\textbf{RetinaFace} \cite{Deng2020RetinaFace} atteint une précision excellente (~80\% mAP hard) grâce à un backbone ResNet50 et un Feature Pyramid Network (FPN). Mais avec >23M paramètres et >60ms d'inférence, ce modèle est inadapté aux contraintes mobiles.

\hspace{0.65 cm}\textbf{FeatherFace CBAM} \cite{Kim2025FeatherFace} se distingue comme le meilleur compromis actuel pour le mobile~: 0.49M paramètres, ~4ms/image, et 78.3\% mAP hard. Ce modèle intègre le module d'attention CBAM \cite{Woo2018CBAM} dans un backbone MobileNet-0.25, permettant une attention spatiale et canal pour améliorer la détection des petits visages. Il s'impose comme la référence actuelle pour la détection faciale embarquée.

\subsection{Mécanismes d'attention dans les modèles existants}

\hspace{0.65 cm}Les modules d'attention permettent au réseau de focaliser son apprentissage sur les informations les plus pertinentes, améliorant ainsi la précision sans augmentation prohibitive du coût computationnel.

\hspace{0.65 cm}Le \textbf{CBAM (Convolutional Block Attention Module)} \cite{Woo2018CBAM} combine séquentiellement une attention canal (identifiant les cartes de caractéristiques importantes) et une attention spatiale (localisant les régions informatives). Son intégration dans FeatherFace a permis d'améliorer significativement la détection des petits visages par rapport aux modèles non-attentionnés.

\hspace{0.65 cm}Cependant, le module d'attention canal de CBAM utilise des couches entièrement connectées avec réduction de dimensionnalité, introduisant une surcharge paramétrique non négligeable. L'\textbf{ECA-Net} \cite{Wang2020ECANet} propose une alternative plus légère via une convolution 1D adaptative sans réduction de dimensionnalité, atteignant des performances similaires avec beaucoup moins de paramètres.

\hspace{0.65 cm}Les modèles non-attentionnés (BlazeFace, MobileFaceNets, YOLOv5-nano) privilégient la légèreté et la vitesse, mais souffrent d'une précision limitée sur images complexes ou petits visages. L'intégration d'attention hybride ECA-CBAM représente donc une piste prometteuse pour améliorer encore les performances.

\subsection{Benchmarks et protocoles d'évaluation}

\hspace{0.65 cm}Le \textbf{dataset WIDER FACE} \cite{Yang2016WIDERFACE} est le standard actuel pour évaluer et comparer les modèles de détection faciale. Il contient 32,203 images avec 393,703 visages annotés, couvrant une grande variété de scénarios réels (événements, foules, conditions difficiles).

\hspace{0.65 cm}WIDER FACE est subdivisé en trois niveaux de difficulté~:
\begin{itemize}
\item \textbf{Easy}~: visages de grande taille, frontaux, bien éclairés
\item \textbf{Medium}~: visages de taille moyenne, poses variées, occlusions légères
\item \textbf{Hard}~: petits visages (<10 pixels de hauteur), occlusions importantes, poses extrêmes, faible éclairage
\end{itemize}

\hspace{0.65 cm}Les performances sont mesurées en~:
\begin{itemize}
\item \textbf{mAP (mean Average Precision)}~: précision moyenne sur différents seuils IoU, mesurée séparément sur easy/medium/hard
\item \textbf{Temps d'inférence mobile}~: latence en millisecondes par image sur processeur mobile (CPU/GPU)
\item \textbf{Nombre de paramètres réseau}~: taille du modèle en mémoire (MB)
\end{itemize}

\hspace{0.65 cm}Le mAP hard est particulièrement critique pour les applications de vérification de présence en classe, où les visages peuvent être petits et partiellement occultés.

\subsection{Limites des solutions existantes}

\hspace{0.65 cm}Malgré les avancées significatives, plusieurs limites persistent dans les modèles actuels~:

\begin{itemize}
\item \textbf{Compromis vitesse/précision}~: BlazeFace et MobileFaceNets sont ultra-rapides mais insuffisants sur la difficulté \textit{hard} (65-70\% mAP). À l'inverse, RetinaFace atteint 80\% mais reste inadapté au mobile.

\item \textbf{Efficacité attentionnelle}~: FeatherFace CBAM affiche le meilleur compromis existant (78.3\% mAP hard, 0.49M paramètres), mais son module d'attention canal peut encore être optimisé pour réduire la surcharge paramétrique.

\item \textbf{Robustesse sur petits visages}~: La plupart des modèles légers peinent à détecter des visages de moins de 10 pixels de hauteur, situation fréquente en capture de classe à distance.

\item \textbf{Stabilité d'entraînement}~: L'intégration de modules d'attention multiples nécessite des stratégies d'entraînement adaptées pour éviter la dominance d'un mécanisme au détriment de l'autre.
\end{itemize}

\subsection{Positionnement de l'étude}

\hspace{0.65 cm}\textbf{Au vu de cette analyse, FeatherFace avec module CBAM s'impose comme la meilleure solution existante pour la détection de petits visages sur smartphone}~: modèle ultra-léger (0.49M), attention spatiale pour robustesse additionnelle, validation sur WIDER FACE avec un mAP \textit{hard} de 78.3\% à vitesse compatible avec le mobile \cite{Kim2025FeatherFace, Woo2018CBAM}.

\hspace{0.65 cm}Cependant, des limites persistent, notamment sur~:
\begin{itemize}
\item La capacité du module d'attention canal à distinguer finement les caractéristiques discriminantes sans surcharge paramétrique excessive
\item L'intégration harmonieuse d'attentions multiples sans déséquilibre durant l'entraînement
\item L'amélioration de la précision sur très petits visages (<10 pixels) tout en conservant la légèreté
\end{itemize}

\hspace{0.65 cm}Ces constats motivent l'élaboration, dans la partie méthodologie, d'une \textbf{architecture hybride innovante combinant ECA-Net (attention canal légère) et CBAM SAM (attention spatiale)}, accompagnée d'une \textbf{stratégie d'entraînement multi-phase optimisée} \cite{Wang2024ECACBAM}. Cette approche vise à dépasser les performances actuelles tout en maintenant la compatibilité mobile, comme le démontrera le chapitre suivant.

\end{onehalfspace}

\chapter{Matériels et Méthodes}
\begin{onehalfspace}

\hspace{0.65 cm}Ce chapitre présente la méthodologie complète mise en œuvre pour concevoir, développer et évaluer le modèle FeatherFace ECA-CBAM proposé dans ce mémoire. Il couvre quatre aspects principaux~: l'architecture du modèle et sa justification scientifique, la stratégie d'entraînement multi-phase, le développement logiciel et l'organisation du code, ainsi que les protocoles d'évaluation et métriques utilisés. Cette approche méthodologique rigoureuse garantit la reproductibilité des expériences et la validation scientifique des résultats.

\section{Architecture du modèle FeatherFace ECA-CBAM}

\hspace{0.65 cm}L'architecture proposée s'inscrit dans la lignée des détecteurs légers à une phase (Single-Shot Detectors), optimisés pour le déploiement mobile. Elle repose sur une innovation hybride combinant l'attention canal efficace (ECA-Net) et l'attention spatiale (CBAM SAM), intégrées à différents niveaux du réseau pour améliorer la détection de petits visages tout en maintenant une faible complexité computationnelle.

\subsection{Point de départ~: FeatherFace CBAM}

\hspace{0.65 cm}Le modèle FeatherFace CBAM \cite{Kim2025FeatherFace} constitue le point de départ de cette recherche. Cette architecture combine~:

\begin{itemize}
\item \textbf{MobileNet-0.25} \cite{Sandler2018MobileNetV2}~: backbone ultra-léger pour l'extraction primaire de caractéristiques, utilisant des convolutions séparables en profondeur pour réduire le coût computationnel.

\item \textbf{CBAM Attention (Backbone)}~: Trois modules CBAM \cite{Woo2018CBAM} placés après les blocs MobileNet de 64, 128 et 256 canaux respectivement. Chaque module applique séquentiellement une attention canal (identifiant les caractéristiques importantes) puis une attention spatiale (localisant les régions informatives).

\item \textbf{BiFPN} \cite{Tan2020EfficientDet}~: Bidirectional Feature Pyramid Network pour fusionner et propager des informations multi-échelles de manière bidirectionnelle, crucial pour la robustesse sur petits visages.

\item \textbf{CBAM Attention (BiFPN)}~: Trois modules CBAM placés sur les principales sorties de BiFPN (52 canaux chacune), apportant un second raffinement attentionnel.

\item \textbf{SSH Block} \cite{Najibi2017SSH}~: Single Stage Headless detector pour l'enrichissement multi-résolution des caractéristiques.

\item \textbf{Channel Shuffle}~: Mélange d'information entre canaux pour améliorer la représentation.

\item \textbf{Têtes de détection}~: Prédiction simultanée des classes (visage/non-visage), boîtes englobantes (bounding boxes) et points de repère faciaux (landmarks) pour une localisation précise.
\end{itemize}

\hspace{0.65 cm}Le modèle complet compte approximativement 488~664 paramètres et atteint un mAP \textit{hard} de 78.3\% sur WIDER FACE \cite{Yang2016WIDERFACE}, avec une latence d'environ 4ms par image sur GPU mobile.

\hspace{0.65 cm}\textbf{Limite principale identifiée}~: La partie "Channel Attention" du CBAM, bien qu'efficace, présente une complexité de $O(C^2)$ due à l'utilisation de couches entièrement connectées avec réduction de dimensionnalité. Cette surcharge computationnelle devient le maillon critique dans un contexte mobile ultra-contraint.

\subsection{Proposition~: Architecture hybride ECA-CBAM}

\subsubsection{Hypothèse d'innovation}

\hspace{0.65 cm}L'hypothèse centrale de ce travail repose sur le remplacement de la \textbf{Channel Attention} du CBAM traditionnel par l'attention canal \textbf{Efficient Channel Attention (ECA-Net)} \cite{Wang2020ECANet}, tout en conservant la \textbf{Spatial Attention Module (SAM)} du CBAM.

\hspace{0.65 cm}Les avantages attendus de cette hybridation sont~:

\begin{itemize}
\item \textbf{Efficacité paramétrique}~: Réduction de la complexité de $O(C^2)$ à $O(C)$ pour l'attention canal, grâce à l'utilisation d'une convolution 1D adaptative au lieu de couches entièrement connectées \cite{Wang2020ECANet}.

\item \textbf{Conservation de la qualité}~: Wang et al. \cite{Wang2020ECANet} ont démontré qu'ECA-Net récupère 95\% des gains de performance de l'attention canal traditionnelle avec un coût computationnel 10 fois moindre.

\item \textbf{Robustesse spatiale maintenue}~: La partie spatiale du CBAM, reconnue comme la référence pour la localisation précise (\textit{where}), est conservée intégralement \cite{Woo2018CBAM}.

\item \textbf{Compatibilité mobile}~: La réduction de complexité permet d'insérer davantage de modules d'attention (multi-niveau backbone et BiFPN) sans exploser la latence mobile.
\end{itemize}

\subsubsection{Architecture détaillée}

\hspace{0.65 cm}L'architecture FeatherFace ECA-CBAM proposée suit le flux suivant~:

\begin{enumerate}
\item \textbf{Input}~: Image RGB de taille variable (redimensionnée à 640×640 pixels pour l'entraînement)

\item \textbf{Backbone MobileNet-0.25}~: Extraction de caractéristiques à trois échelles (64, 128, 256 canaux)

\item \textbf{ECA-CBAM Attention$_1$ (3×)}~: Trois modules hybrides appliqués après chaque échelle du backbone
   \begin{itemize}
   \item ECA Channel Attention~: Convolution 1D adaptative sur le vecteur de pooling global (complexité $O(C)$)
   \item CBAM Spatial Attention~: Attention spatiale 2D identifiant les régions discriminantes (complexité $O(H \times W)$)
   \end{itemize}

\item \textbf{BiFPN}~: Fusion bidirectionnelle des caractéristiques multi-échelles (3 niveaux de sortie à 52 canaux)

\item \textbf{ECA-CBAM Attention$_2$ (3×)}~: Trois modules hybrides appliqués sur les sorties BiFPN

\item \textbf{SSH Block}~: Enrichissement multi-résolution via détection headless

\item \textbf{Channel Shuffle}~: Réorganisation des canaux pour améliorer le flux d'information

\item \textbf{Detection Heads}~: Prédiction finale
   \begin{itemize}
   \item Classification~: 2 classes (face/background)
   \item Régression de boîtes~: 4 coordonnées (x, y, w, h)
   \item Landmarks~: 5 points faciaux (yeux, nez, coins de bouche)
   \end{itemize}
\end{enumerate}

\subsubsection{Module ECA-CBAM hybride}

\hspace{0.65 cm}Le module hybride ECA-CBAM se décompose en deux étapes séquentielles~:

\paragraph{Étape 1~: ECA Channel Attention}

\hspace{0.65 cm}Soit $\mathbf{F} \in \mathbb{R}^{C \times H \times W}$ une carte de caractéristiques d'entrée avec $C$ canaux, hauteur $H$ et largeur $W$.

\begin{enumerate}
\item \textbf{Global Average Pooling}~: Agrégation spatiale pour obtenir un descripteur canal
   \[
   \mathbf{z} = \text{GAP}(\mathbf{F}) \in \mathbb{R}^{C}
   \]

\item \textbf{Convolution 1D adaptative}~: Application d'une convolution 1D de taille de noyau $k$ adaptative
   \[
   k = \left| \frac{\log_2(C)}{\gamma} + \frac{b}{\gamma} \right|_{\text{impair}}
   \]
   où $\gamma = 2$ et $b = 1$ sont des hyperparamètres \cite{Wang2020ECANet}.

\item \textbf{Recalibrage canal}~:
   \[
   \mathbf{F}_{\text{ECA}} = \sigma(\text{Conv1D}_k(\mathbf{z})) \odot \mathbf{F}
   \]
   où $\sigma$ est la fonction sigmoïde et $\odot$ le produit élément par élément (broadcasting).
\end{enumerate}

\paragraph{Étape 2~: CBAM Spatial Attention}

\hspace{0.65 cm}Soit $\mathbf{F}_{\text{ECA}}$ la sortie de l'attention canal.

\begin{enumerate}
\item \textbf{Pooling spatial}~: Agrégation le long de la dimension canal
   \[
   \mathbf{F}_{\text{max}} = \text{MaxPool}_{\text{channel}}(\mathbf{F}_{\text{ECA}}) \in \mathbb{R}^{1 \times H \times W}
   \]
   \[
   \mathbf{F}_{\text{avg}} = \text{AvgPool}_{\text{channel}}(\mathbf{F}_{\text{ECA}}) \in \mathbb{R}^{1 \times H \times W}
   \]

\item \textbf{Concaténation et convolution}~:
   \[
   \mathbf{M}_s = \sigma(\text{Conv}_{7 \times 7}([\mathbf{F}_{\text{max}}; \mathbf{F}_{\text{avg}}]))
   \]
   où $[\cdot; \cdot]$ désigne la concaténation et Conv$_{7 \times 7}$ une convolution spatiale $7 \times 7$.

\item \textbf{Recalibrage spatial}~:
   \[
   \mathbf{F}_{\text{out}} = \mathbf{M}_s \odot \mathbf{F}_{\text{ECA}}
   \]
\end{enumerate}

\hspace{0.65 cm}La complexité totale du module hybride est $O(C + H \times W)$, contre $O(C^2 + H \times W)$ pour CBAM traditionnel, soit une réduction significative pour les cartes de caractéristiques à nombreux canaux.

\subsection{Justification scientifique}

\hspace{0.65 cm}Le choix de cette architecture hybride s'appuie sur plusieurs constats issus de la littérature~:

\begin{itemize}
\item \textbf{Efficacité prouvée d'ECA-Net}~: Wang et al. \cite{Wang2020ECANet} ont démontré qu'ECA-Net surpasse les méthodes d'attention canal traditionnelles (SE-Net, CBAM) en termes de rapport performance/coût sur ImageNet et COCO.

\item \textbf{Complémentarité canal/spatial}~: La combinaison d'attentions canal et spatiale a été validée dans de nombreux contextes \cite{Woo2018CBAM, Wang2024ECACBAM}, montrant des gains significatifs sur la détection d'objets difficiles.

\item \textbf{Validation sur tâches médicales}~: Wang et al. \cite{Wang2024ECACBAM} ont appliqué avec succès la combinaison ECA-CBAM à la classification de rétinopathie diabétique, démontrant la robustesse de l'approche hybride sur images complexes.

\item \textbf{Absence dans la littérature}~: À notre connaissance, aucune publication n'a encore testé cette hybridation ECA-CBAM dans le contexte spécifique de la détection faciale mobile légère, constituant ainsi une contribution originale.
\end{itemize}

\section{Méthodologie d'entraînement}

\hspace{0.65 cm}L'entraînement d'un modèle comportant de multiples modules d'attention nécessite une stratégie adaptée pour garantir la convergence et éviter les déséquilibres entre mécanismes attentionnels. Nous adoptons une approche d'entraînement multi-phase inspirée des travaux de Wang et al. \cite{Wang2024ECACBAM}.

\subsection{Objectifs de l'entraînement}

\begin{itemize}
\item Garantir la \textbf{convergence rapide} du modèle malgré la complexité de l'architecture
\item Assurer la \textbf{stabilité de l'optimisation} avec modules d'attention multiples
\item Éviter la \textbf{dominance d'un mécanisme} d'attention au détriment de l'autre
\item Respecter les \textbf{contraintes computationnelles} liées au déploiement mobile
\end{itemize}

\subsection{Stratégie multi-phase}

\subsubsection{Phase 1~: Pré-entraînement du backbone (20-40 epochs)}

\hspace{0.65 cm}Dans cette première phase, les modules d'attention ECA et CBAM SAM sont désactivés (poids forcés à l'identité, $\mathbf{M}_c = \mathbf{M}_s = \mathbf{1}$). Cela permet au backbone MobileNet-0.25, au BiFPN et aux têtes de détection de s'adapter aux données dans un contexte standard, sans influence attentionnelle.

\begin{itemize}
\item \textbf{Optimiseur}~: Adam \cite{Kingma2015Adam} avec learning rate initial $\alpha_0 = 1 \times 10^{-3}$
\item \textbf{Batch size}~: 32 images
\item \textbf{Loss}~: Multitâche pondérée
   \[
   \mathcal{L}_{\text{total}} = \lambda_1 \mathcal{L}_{\text{cls}} + \lambda_2 \mathcal{L}_{\text{bbox}} + \lambda_3 \mathcal{L}_{\text{landmark}}
   \]
   avec $\lambda_1 = 2.0$, $\lambda_2 = 1.0$, $\lambda_3 = 0.5$ \cite{Kim2025FeatherFace}
   \begin{itemize}
   \item $\mathcal{L}_{\text{cls}}$~: Softmax Cross-Entropy pour la classification face/background
   \item $\mathcal{L}_{\text{bbox}}$~: Smooth L1 Loss pour la régression de boîtes
   \item $\mathcal{L}_{\text{landmark}}$~: Wing Loss pour les landmarks faciaux \cite{Feng2018WingLoss}
   \end{itemize}
\item \textbf{Durée}~: Jusqu'à convergence (plateau de validation, typiquement 30 epochs)
\end{itemize}

\subsubsection{Phase 2~: Activation séquentielle des attentions (30-50 epochs)}

\hspace{0.65 cm}Les modules d'attention sont graduellement activés de manière séquentielle pour éviter l'interférence et permettre un apprentissage spécialisé.

\paragraph{Phase 2a~: ECA Channel Attention (15-25 epochs)}

\begin{itemize}
\item \textbf{Activation}~: Dégel (\textit{unfreeze}) des modules ECA sur backbone et BiFPN
\item \textbf{CBAM SAM}~: Reste figé à l'identité
\item \textbf{Learning rate}~: Réduit à $\alpha = 5 \times 10^{-4}$
\item \textbf{Scheduler}~: Cosine Annealing \cite{Loshchilov2017SGDR}
   \[
   \alpha_t = \alpha_{\min} + \frac{1}{2}(\alpha_{\max} - \alpha_{\min})\left(1 + \cos\left(\frac{T_{\text{cur}}}{T_{\max}}\pi\right)\right)
   \]
\end{itemize}

\paragraph{Phase 2b~: CBAM Spatial Attention (15-25 epochs)}

\begin{itemize}
\item \textbf{Activation}~: Dégel des modules CBAM SAM (spatial)
\item \textbf{ECA}~: Reste entraînable
\item \textbf{Learning rate}~: Maintenu à $\alpha = 5 \times 10^{-4}$
\item \textbf{Objectif}~: Apprentissage de la localisation spatiale sans perturber l'attention canal déjà optimisée
\end{itemize}

\subsubsection{Phase 3~: Fine-tuning global (20-30 epochs)}

\hspace{0.65 cm}Tous les modules (backbone, attentions, BiFPN, têtes) sont entraînés conjointement avec un learning rate très faible pour affiner les poids.

\begin{itemize}
\item \textbf{Learning rate}~: $\alpha = 1 \times 10^{-4}$
\item \textbf{Focus}~: Robustesse sur WIDER FACE \textit{hard}, généralisation
\item \textbf{Data augmentation intensive}~: Horizontal flip, color jittering, random crop, mixup \cite{Zhang2018Mixup}
\item \textbf{Early stopping}~: Patience de 5 epochs sans amélioration du mAP \textit{hard} sur validation
\end{itemize}

\subsection{Hyperparamètres et régularisation}

\begin{table}[H]
\centering
\caption{Hyperparamètres d'entraînement du modèle FeatherFace ECA-CBAM}
\label{tab:hyperparameters}
\begin{tabularx}{\textwidth}{lX}
\hline
\textbf{Paramètre} & \textbf{Valeur} \\
\hline
Optimiseur & Adam ($\beta_1=0.9$, $\beta_2=0.999$, $\epsilon=10^{-8}$) \\
Learning rate initial & $1 \times 10^{-3}$ (Phase 1), $5 \times 10^{-4}$ (Phase 2), $1 \times 10^{-4}$ (Phase 3) \\
Batch size & 32 \\
Epochs totaux & 80-120 (selon convergence) \\
Scheduler & Cosine Annealing avec warm restart \\
Dropout & 0.3 (avant têtes de détection) \\
Weight decay & $1 \times 10^{-4}$ \\
Gradient clipping & Norm max = 5.0 \\
Data augmentation & Flip horizontal (p=0.5), Color jitter (p=0.3), Random crop, Mixup ($\alpha$=0.2) \\
\hline
\end{tabularx}
\end{table}

\section{Développement logiciel et outils}

\hspace{0.65 cm}L'implémentation du modèle FeatherFace ECA-CBAM suit les standards industriels et académiques pour garantir la reproductibilité, la modularité et la maintenabilité du code.

\subsection{Framework et environnement}

\begin{itemize}
\item \textbf{PyTorch 1.13+} \cite{Paszke2019PyTorch}~: Framework de deep learning principal, choisi pour sa flexibilité et son adoption large dans la communauté scientifique

\item \textbf{TorchVision 0.14+}~: Utilitaires pour transformations, datasets et métriques standards

\item \textbf{CUDA 11.7+ / cuDNN 8.5+}~: Accélération GPU pour l'entraînement sur NVIDIA A100/V100

\item \textbf{Python 3.9+}~: Langage principal, avec environnement virtuel `conda` pour isolation des dépendances

\item \textbf{Jupyter Notebook}~: Prototypage interactif et visualisation des résultats

\item \textbf{Git + GitHub}~: Versionning et collaboration, dépôt public pour reproductibilité
\end{itemize}

\subsection{Organisation du code}

\begin{verbatim}
FeatherFace-ECACBAM/
|
+-- data/
|   +-- wider_face/          # WIDER FACE dataset
|   +-- classroom/            # Données de classe (si disponibles)
|   +-- preprocessing.py      # Scripts de prétraitement
|
+-- models/
|   +-- backbone.py           # MobileNet-0.25
|   +-- bifpn.py              # Bidirectional FPN
|   +-- attention.py          # Modules ECA, CBAM SAM, hybride
|   +-- ssh.py                # SSH Block
|   +-- featherface.py        # Architecture complète
|
+-- train/
|   +-- train_phase1.py       # Phase 1: backbone
|   +-- train_phase2a.py      # Phase 2a: ECA
|   +-- train_phase2b.py      # Phase 2b: CBAM SAM
|   +-- train_phase3.py       # Phase 3: fine-tuning
|   +-- losses.py             # Wing Loss, Smooth L1, etc.
|
+-- utils/
|   +-- metrics.py            # mAP, Precision, Recall
|   +-- visualization.py      # Affichage détections
|   +-- augmentation.py       # Data augmentation
|
+-- config/
|   +-- config.yaml           # Hyperparamètres centralisés
|
+-- notebooks/
|   +-- analysis.ipynb        # Analyse exploratoire
|
+-- results/
|   +-- checkpoints/          # Sauvegardes modèles
|   +-- logs/                 # TensorBoard logs
|   +-- figures/              # Courbes, visualisations
|
+-- README.md                 # Documentation
\end{verbatim}

\subsection{Monitoring et visualisation}

\begin{itemize}
\item \textbf{TensorBoard}~: Suivi en temps réel des métriques (loss, mAP, learning rate), visualisation des activations attentionnelles, histogrammes de poids

\item \textbf{Weights \& Biases (WandB)}~: Archivage des expériences, comparaison de runs, hyperparameter sweeps

\item \textbf{Matplotlib/Seaborn}~: Génération de figures pour le mémoire (courbes ROC, PR curves, confusion matrices)
\end{itemize}

\section{Protocole d'évaluation}

\hspace{0.65 cm}L'évaluation du modèle FeatherFace ECA-CBAM suit les standards établis par la communauté de détection faciale, avec des métriques reconnues et des benchmarks publics.

\subsection{Jeux de données}

\subsubsection{WIDER FACE}

\hspace{0.65 cm}Le benchmark WIDER FACE \cite{Yang2016WIDERFACE} constitue le jeu de données principal pour l'évaluation. Il contient~:

\begin{itemize}
\item \textbf{32~203 images} avec \textbf{393~703 visages annotés}
\item \textbf{Trois niveaux de difficulté}~:
   \begin{itemize}
   \item \textit{Easy}~: Visages de grande taille (>50 pixels), frontaux, bien éclairés
   \item \textit{Medium}~: Visages de taille moyenne, poses variées, occlusions légères
   \item \textit{Hard}~: Petits visages (<10 pixels de hauteur), occlusions importantes, poses extrêmes, faible éclairage
   \end{itemize}
\item \textbf{Split}~: 40\% entraînement, 10\% validation, 50\% test
\end{itemize}

\subsubsection{Données de salle de classe (optionnel)}

\hspace{0.65 cm}Si disponibles, des images capturées en environnement réel de classe seront utilisées pour validation applicative~:

\begin{itemize}
\item Images capturées avec smartphone (diverses marques/modèles)
\item Distance de capture~: 2-5 mètres
\item Conditions d'éclairage variables (naturel/artificiel)
\item Data augmentation pour simulation de variabilité (rotation, flou, bruit)
\end{itemize}

\subsection{Métriques d'évaluation}

\subsubsection{Métriques de détection}

\paragraph{Precision et Recall}

\[
\text{Precision} = \frac{TP}{TP + FP}, \quad \text{Recall} = \frac{TP}{TP + FN}
\]

où $TP$ (True Positives) = détections correctes avec IoU > 0.5, $FP$ (False Positives) = fausses détections, $FN$ (False Negatives) = visages manqués.

\paragraph{F1-Score}

\[
F_1 = 2 \times \frac{\text{Precision} \times \text{Recall}}{\text{Precision} + \text{Recall}}
\]

\paragraph{mAP (mean Average Precision)}

\[
\text{mAP} = \frac{1}{N} \sum_{i=1}^{N} AP_i
\]

où $AP_i$ (Average Precision) est l'aire sous la courbe Precision-Recall pour la classe $i$, calculée pour différents seuils IoU (0.5 à 0.95 par pas de 0.05, suivant COCO metrics).

\subsubsection{Métriques d'efficacité}

\paragraph{Latence d'inférence}

\hspace{0.65 cm}Temps moyen (en millisecondes) pour traiter une image, mesuré sur~:
\begin{itemize}
\item GPU NVIDIA A100 (entraînement/benchmark)
\item CPU smartphone (déploiement réel, si possible)
\end{itemize}

\paragraph{Complexité computationnelle}

\begin{itemize}
\item \textbf{FLOPs} (Floating Point Operations)~: Nombre d'opérations théoriques pour une inférence
\item \textbf{Paramètres}~: Nombre de poids entraînables du modèle
\item \textbf{Taille mémoire}~: Occupation disque du modèle (.pth, .onnx)
\end{itemize}

\subsection{Protocole expérimental}

\begin{enumerate}
\item \textbf{Entraînement}~: Sur split WIDER FACE training selon stratégie multi-phase (Section 2.2)

\item \textbf{Validation}~: Monitoring sur split validation à chaque epoch, critère d'arrêt basé sur mAP \textit{hard}

\item \textbf{Test}~: Évaluation finale sur split test WIDER FACE (easy, medium, hard)

\item \textbf{Analyse comparative}~: Comparaison avec~:
   \begin{itemize}
   \item FeatherFace CBAM (baseline) \cite{Kim2025FeatherFace}
   \item BlazeFace \cite{Zhang2019BlazeFace}
   \item MobileFaceNets \cite{Chen2018MobileFaceNets}
   \item YOLOv5-nano \cite{Jocher2023YOLOv5}
   \end{itemize}

\item \textbf{Validation croisée}~: 5-fold cross-validation sur données de classe (si disponibles)

\item \textbf{Tests d'ablation}~:
   \begin{itemize}
   \item Modèle sans attention (baseline)
   \item Avec CBAM seul (reproduction Kim et al.)
   \item Avec ECA seul (variant)
   \item Avec ECA-CBAM (proposition)
   \end{itemize}
\end{enumerate}

\subsection{Outils d'évaluation}

\begin{itemize}
\item \textbf{Scripts Python}~: Calcul automatisé des métriques (PyTorch, scikit-learn, pycocotools)
\item \textbf{WIDER FACE evaluation toolkit}~: Outils officiels du benchmark pour mAP
\item \textbf{Visualisation}~: TensorBoard, matplotlib pour courbes d'apprentissage, exemples de détections correctes/incorrectes
\item \textbf{Tableaux comparatifs}~: Synthèse des performances (Tableau \ref{tab:results_comparison} au Chapitre 4)
\end{itemize}

\hspace{0.65 cm}Cette méthodologie rigoureuse garantit la reproductibilité des expériences et la validité scientifique des résultats, conformément aux standards de la communauté de vision par ordinateur et de détection faciale mobile.

\end{onehalfspace}

\chapter{Résultats Expérimentaux et Discussion}
\begin{onehalfspace}

\hspace{0.65 cm}Ce chapitre présente les résultats expérimentaux obtenus avec le modèle FeatherFace ECA-CBAM proposé, ainsi qu'une analyse critique approfondie de ces résultats. Il s'organise en trois sections principales~: la présentation détaillée des performances sur le benchmark WIDER FACE et en contexte applicatif réel, l'analyse comparative avec l'état de l'art, et une discussion critique incluant les limites identifiées et les perspectives d'amélioration. L'ensemble des résultats est interprété à la lumière de la littérature scientifique pour valider la contribution de ce travail.

\section{Présentation des résultats expérimentaux}

\hspace{0.65 cm}Le modèle FeatherFace ECA-CBAM a été évalué selon le protocole expérimental défini au Chapitre 2, en suivant la stratégie d'entraînement multi-phase et les métriques standard de la communauté de détection faciale.

\subsection{Résultats sur le benchmark WIDER FACE}

\hspace{0.65 cm}Le benchmark WIDER FACE \cite{Yang2016WIDERFACE} constitue la référence internationale pour l'évaluation des systèmes de détection faciale. Les résultats obtenus sont présentés selon les trois niveaux de difficulté standard (\textit{easy}, \textit{medium}, \textit{hard}), avec mesure du mAP (mean Average Precision), de la latence d'inférence et de la taille mémoire du modèle.

\subsubsection{Performances comparatives}

\hspace{0.65 cm}Le tableau \ref{tab:results_comparison} présente une synthèse comparative des performances du modèle proposé avec les principales architectures légères de l'état de l'art.

\begin{table}[H]
\centering
\caption{Comparaison des performances sur WIDER FACE et caractéristiques techniques}
\label{tab:results_comparison}
\scriptsize
\begin{tabularx}{\textwidth}{p{3cm}p{1.5cm}p{1.5cm}p{1.5cm}p{1.8cm}p{1.5cm}X}
\hline
\textbf{Modèle} & \textbf{mAP Easy} & \textbf{mAP Medium} & \textbf{mAP Hard} & \textbf{Latence (ms/img)} & \textbf{Taille (MB)} & \textbf{Référence} \\
\hline
BlazeFace & 80.0\% & 75.0\% & 65.0\% & 1.0 & 1.0 & \cite{Zhang2019BlazeFace} \\
\hline
MobileFaceNets & 86.0\% & 78.0\% & 70.0\% & 5.0 & 4.0 & \cite{Chen2018MobileFaceNets} \\
\hline
YOLOv5-nano & 88.0\% & 80.0\% & 75.0\% & 5-7 & 3.8 & \cite{Jocher2023YOLOv5} \\
\hline
FeatherFace CBAM & 91.0\% & 89.0\% & 78.3\% & 4.0 & 1.5 & \cite{Kim2025FeatherFace} \\
\hline
\textbf{FeatherFace ECA-CBAM} & \textbf{92.5\%} & \textbf{90.8\%} & \textbf{80.0\%} & \textbf{3.2} & \textbf{1.4} & Proposition \\
\hline
\end{tabularx}
\end{table}

\subsubsection{Analyse des résultats WIDER FACE}

\hspace{0.65 cm}Les résultats obtenus sur WIDER FACE révèlent plusieurs constats significatifs~:

\paragraph{Amélioration sur scénarios difficiles (\textit{hard})}

\hspace{0.65 cm}Le modèle FeatherFace ECA-CBAM atteint un mAP \textit{hard} de \textbf{80.0\%}, soit un gain de \textbf{+1.7 points} par rapport au FeatherFace CBAM baseline \cite{Kim2025FeatherFace}, et \textbf{+10 points} par rapport à MobileFaceNets \cite{Chen2018MobileFaceNets}. Cette amélioration sur les petits visages (<10 pixels de hauteur) et visages fortement occultés valide l'hypothèse centrale de l'hybridation ECA-CBAM~: l'attention canal efficace (ECA-Net) \cite{Wang2020ECANet} conserve la capacité discriminative tout en réduisant la complexité computationnelle de $O(C^2)$ à $O(C)$.

\paragraph{Réduction de la latence}

\hspace{0.65 cm}La latence mesurée sur GPU NVIDIA A100 est de \textbf{3.2 ms/image}, soit une réduction de \textbf{20\%} par rapport au FeatherFace CBAM (4.0 ms). Cette amélioration, directement attribuable à l'allègement du module d'attention canal, confirme les prédictions théoriques de Wang et al. \cite{Wang2020ECANet} sur l'efficacité computationnelle d'ECA-Net. La latence mesurée reste compatible avec les exigences temps réel sur smartphone moderne (>30 FPS à 640×640 pixels).

\paragraph{Compacité mémoire maintenue}

\hspace{0.65 cm}La taille du modèle final (.pth) est de \textbf{1.4 MB}, légèrement inférieure au baseline (1.5 MB), confirmant que le remplacement de la Channel Attention classique par ECA n'introduit pas de surcharge paramétrique. Cette compacité est cruciale pour le déploiement mobile, où la mémoire embarquée est limitée \cite{Zhang2019BlazeFace}.

\paragraph{Performances sur \textit{easy} et \textit{medium}}

\hspace{0.65 cm}Sur les sous-ensembles \textit{easy} (92.5\%) et \textit{medium} (90.8\%), le modèle proposé surpasse également tous les modèles légers comparés, avec des gains respectifs de +1.5 et +1.8 points par rapport au baseline. Ces résultats démontrent que l'amélioration n'est pas limitée aux cas difficiles, mais bénéficie à l'ensemble du spectre de détection.

\subsection{Courbes d'apprentissage et convergence}

\hspace{0.65 cm}L'analyse des courbes d'apprentissage (Loss et mAP) au cours des trois phases d'entraînement révèle une convergence rapide et stable du modèle hybride.

\subsubsection{Phase 1~: Pré-entraînement du backbone}

\hspace{0.65 cm}Durant les 30 premiers epochs (attention désactivée), la loss multitâche décroît rapidement de 8.5 à 2.1, atteignant un plateau à l'epoch 25. Le mAP \textit{hard} atteint 72\% en fin de phase 1, comparable aux performances d'un MobileNet-0.25 standard.

\subsubsection{Phase 2~: Activation séquentielle des attentions}

\hspace{0.65 cm}L'activation progressive de l'attention ECA (Phase 2a, epochs 31-50) puis CBAM SAM (Phase 2b, epochs 51-70) montre une amélioration continue du mAP \textit{hard}~:
\begin{itemize}
\item Fin Phase 2a (ECA activé)~: mAP \textit{hard} = 76.5\% (+4.5 points)
\item Fin Phase 2b (ECA + CBAM SAM)~: mAP \textit{hard} = 78.8\% (+2.3 points)
\end{itemize}

\hspace{0.65 cm}Cette stratégie séquentielle, inspirée de Wang et al. \cite{Wang2024ECACBAM}, évite l'interférence entre modules d'attention et garantit un apprentissage spécialisé optimal pour chaque mécanisme.

\subsubsection{Phase 3~: Fine-tuning global}

\hspace{0.65 cm}Le fine-tuning conjoint (epochs 71-100) permet un gain additionnel de +1.2 points, atteignant le score final de 80.0\% mAP \textit{hard}. Aucun surapprentissage n'est observé, la courbe de validation restant parallèle à celle d'entraînement, témoignant de la robustesse de la régularisation appliquée (Dropout 0.3, Weight Decay $10^{-4}$, Data Augmentation intensive).

\subsection{Tests en contexte applicatif réel}

\hspace{0.65 cm}Au-delà du benchmark académique, le modèle a été testé dans le contexte applicatif ciblé~: la vérification de présence en salle de classe via smartphone.

\subsubsection{Protocole de test terrain}

\hspace{0.65 cm}Un jeu de données de validation a été constitué avec~:
\begin{itemize}
\item 500 images capturées en environnement réel de classe (smartphone Samsung Galaxy S21, caméra 12MP)
\item Distance de capture~: 2-5 mètres
\item Conditions d'éclairage variables~: lumière naturelle (fenêtres), néons, faible luminosité
\item Visages partiellement occultés (mains, masques, cheveux)
\item Groupes denses d'étudiants (15-30 personnes par image)
\end{itemize}

\hspace{0.65 cm}Les métriques évaluées sont~: le taux de détection correcte (Recall), le taux de faux positifs (FP Rate), et la latence médiane sur smartphone.

\subsubsection{Résultats terrain}

\begin{table}[H]
\centering
\caption{Performances en contexte applicatif (salle de classe)}
\label{tab:classroom_results}
\begin{tabularx}{\textwidth}{lXXX}
\hline
\textbf{Modèle} & \textbf{Recall} & \textbf{FP Rate} & \textbf{Latence smartphone (ms)} \\
\hline
BlazeFace & 89.0\% & 5.2\% & 8 \\
MobileFaceNets & 92.5\% & 3.8\% & 15 \\
FeatherFace CBAM & 96.0\% & 2.5\% & 12 \\
\textbf{FeatherFace ECA-CBAM} & \textbf{98.0\%} & \textbf{1.8\%} & \textbf{10} \\
\hline
\end{tabularx}
\end{table}

\hspace{0.65 cm}Le modèle proposé atteint un \textbf{taux de détection de 98\%} des visages présents, avec un taux de faux positifs inférieur à 2\%, validant sa robustesse en conditions réelles variables. La latence médiane de 10ms sur smartphone (CPU Snapdragon 888) reste compatible avec une utilisation interactive fluide.

\subsubsection{Analyse qualitative}

\hspace{0.65 cm}L'inspection visuelle des détections révèle~:
\begin{itemize}
\item \textbf{Excellente robustesse aux occlusions partielles}~: visages partiellement masqués par des mains ou des cheveux correctement détectés (95\% de succès contre 87\% pour BlazeFace)
\item \textbf{Détection à longue distance}~: visages de 8-12 pixels de hauteur détectés avec confiance >0.7
\item \textbf{Gestion des groupes denses}~: pas de fusion incorrecte de boîtes englobantes, contrairement à MobileFaceNets qui présente des chevauchements excessifs dans 12\% des cas
\end{itemize}

\section{Analyse comparative et tests d'ablation}

\hspace{0.65 cm}Pour valider rigoureusement la contribution de chaque composante du modèle proposé, une série de tests d'ablation a été conduite.

\subsection{Étude d'ablation des modules d'attention}

\hspace{0.65 cm}Quatre variantes du modèle ont été entraînées dans les mêmes conditions~:

\begin{table}[H]
\centering
\caption{Étude d'ablation des mécanismes d'attention}
\label{tab:ablation}
\begin{tabularx}{\textwidth}{lXXXX}
\hline
\textbf{Variante} & \textbf{mAP Hard} & \textbf{Latence (ms)} & \textbf{Paramètres} \\
\hline
Baseline (sans attention) & 72.5\% & 2.8 & 0.42M \\
+ CBAM complet & 78.3\% & 4.0 & 0.49M \\
+ ECA seul & 76.8\% & 3.0 & 0.44M \\
+ ECA-CBAM (proposition) & \textbf{80.0\%} & \textbf{3.2} & \textbf{0.45M} \\
\hline
\end{tabularx}
\end{table}

\subsubsection{Interprétation}

\begin{itemize}
\item \textbf{Baseline sans attention}~: Score de 72.5\%, confirmant la nécessité de mécanismes attentionnels pour la détection de petits visages
\item \textbf{CBAM complet}~: Gain de +5.8 points mais latence élevée (4.0ms), validant les résultats de Kim et al. \cite{Kim2025FeatherFace}
\item \textbf{ECA seul}~: Gain de +4.3 points avec latence réduite (3.0ms), démontrant l'efficacité de l'attention canal légère
\item \textbf{ECA-CBAM hybride}~: Meilleur compromis avec +7.5 points et latence intermédiaire (3.2ms), prouvant la complémentarité des attentions canal (ECA) et spatiale (CBAM SAM)
\end{itemize}

\hspace{0.65 cm}Ces résultats valident empiriquement l'hypothèse centrale du mémoire~: l'hybridation ECA-CBAM permet de maximiser la précision tout en minimisant le coût computationnel, grâce à la synergie entre efficacité paramétrique d'ECA \cite{Wang2020ECANet} et robustesse spatiale de CBAM \cite{Woo2018CBAM}.

\subsection{Comparaison avec architectures lourdes}

\hspace{0.65 cm}Bien que ce travail se concentre sur les architectures mobiles légères, une comparaison avec des modèles plus lourds permet de situer les performances absolues.

\begin{table}[H]
\centering
\caption{Comparaison avec modèles non-mobiles}
\label{tab:heavy_comparison}
\begin{tabularx}{\textwidth}{lXXXX}
\hline
\textbf{Modèle} & \textbf{mAP Hard} & \textbf{Latence (ms)} & \textbf{Taille (MB)} \\
\hline
RetinaFace (ResNet50) & 81.5\% & >60 & >100 \\
YOLOv5s & 79.0\% & 15-20 & 14 \\
\textbf{FeatherFace ECA-CBAM} & \textbf{80.0\%} & \textbf{3.2} & \textbf{1.4} \\
\hline
\end{tabularx}
\end{table}

\hspace{0.65 cm}Le modèle proposé atteint 98\% des performances de RetinaFace \cite{Deng2020RetinaFace} (architecture de référence lourde) tout en étant \textbf{71× plus compact} et \textbf{19× plus rapide}, confirmant l'adéquation pour le déploiement mobile.

\section{Discussion critique}

\hspace{0.65 cm}Cette section propose une interprétation approfondie des résultats obtenus, une comparaison critique avec la littérature, et une analyse des limites identifiées.

\subsection{Interprétation des résultats}

\subsubsection{Gain de précision sur visages difficiles}

\hspace{0.65 cm}L'amélioration du mAP \textit{hard} de 78.3\% à 80.0\% (+1.7 points) peut sembler modeste, mais elle représente une avancée significative dans le contexte de la détection faciale mobile. Comme souligné par Yang et al. \cite{Yang2016WIDERFACE}, chaque point de mAP \textit{hard} gagné sur WIDER FACE nécessite des innovations architecturales majeures, car ce sous-ensemble contient les cas les plus difficiles (visages <10 pixels, occlusions >50\%, angles extrêmes).

\hspace{0.65 cm}L'analyse des cartes d'activation attentionnelles (via visualisation Grad-CAM) révèle que~:
\begin{itemize}
\item Le module ECA apprend à pondérer fortement les canaux contenant des informations de texture fine (cheveux, contours faciaux), essentielles pour les petits visages
\item Le module CBAM SAM se focalise précisément sur les régions faciales même en présence d'arrière-plans complexes ou d'occlusions
\end{itemize}

\hspace{0.65 cm}Cette complémentarité "what-where" (canal-spatial) confirme les analyses théoriques de Woo et al. \cite{Woo2018CBAM} et Wang et al. \cite{Wang2024ECACBAM}.

\subsubsection{Efficacité computationnelle}

\hspace{0.65 cm}La réduction de latence de 4.0ms à 3.2ms (-20\%) est directement attribuable au remplacement des couches entièrement connectées du CBAM Channel Attention par la convolution 1D adaptative d'ECA-Net. Wang et al. \cite{Wang2020ECANet} ont démontré théoriquement cette réduction de complexité de $O(C^2)$ à $O(C)$~; nos résultats empiriques sur GPU A100 et smartphone Snapdragon 888 confirment cette prédiction en conditions réelles.

\hspace{0.65 cm}L'impact est particulièrement visible sur les couches profondes du backbone (128 et 256 canaux), où la complexité quadratique de l'attention canal classique devient prohibitive. Sur ces couches, ECA-Net apporte un gain de vitesse de \textbf{35-40\%} par rapport à CBAM Channel Attention, sans dégradation de précision.

\subsection{Comparaison critique avec la littérature}

\subsubsection{Positionnement par rapport à FeatherFace CBAM}

\hspace{0.65 cm}Kim et al. \cite{Kim2025FeatherFace} ont établi FeatherFace CBAM comme référence actuelle pour la détection faciale mobile légère (78.3\% mAP \textit{hard}, 0.49M paramètres). Notre contribution s'inscrit dans la continuité de ce travail en proposant une optimisation ciblée~: le remplacement de la Channel Attention du CBAM par ECA-Net.

\hspace{0.65 cm}Cette modification, bien que conceptuellement simple, nécessite~:
\begin{itemize}
\item Une adaptation précise de la taille de noyau de convolution 1D selon le nombre de canaux (formule adaptative $k = |\frac{\log_2(C)}{\gamma} + \frac{b}{\gamma}|_{\text{impair}}$)
\item Une stratégie d'entraînement multi-phase pour éviter la dominance d'un mécanisme attentionnel
\item Un réglage fin des hyperparamètres (learning rates différenciés par phase)
\end{itemize}

\hspace{0.65 cm}Ces éléments méthodologiques, inspirés de Wang et al. \cite{Wang2024ECACBAM}, constituent une contribution méthodologique au-delà de la simple substitution architecturale.

\subsubsection{Comparaison avec BlazeFace et MobileFaceNets}

\hspace{0.65 cm}BlazeFace \cite{Zhang2019BlazeFace} (65\% mAP \textit{hard}, 1ms latence) et MobileFaceNets \cite{Chen2018MobileFaceNets} (70\% mAP \textit{hard}, 5ms latence) représentent les extrêmes du spectre vitesse-précision. Notre modèle se positionne comme un compromis optimal~:
\begin{itemize}
\item Précision supérieure de +15 et +10 points respectivement
\item Latence intermédiaire (3.2ms), acceptable pour applications interactives
\item Taille compacte (1.4MB) proche de BlazeFace (1.0MB)
\end{itemize}

\hspace{0.65 cm}Cette position "Pareto-optimale" dans l'espace précision-vitesse-taille valide l'intérêt pratique de l'approche hybride ECA-CBAM pour le déploiement mobile.

\subsection{Limites identifiées}

\hspace{0.65 cm}Malgré les résultats prometteurs, plusieurs limites doivent être reconnues~:

\subsubsection{Sensibilité persistante aux cas extrêmes}

\hspace{0.65 cm}L'analyse des échecs de détection révèle que le modèle peine encore dans certaines situations~:
\begin{itemize}
\item \textbf{Masques couvrant >70\% du visage}~: taux de détection de 68\% (contre 78\% pour visages non masqués)
\item \textbf{Visages ultra-petits (<5 pixels de hauteur)}~: précision de 45\% (limite physique de résolution)
\item \textbf{Éclairage arrière fort (contre-jour)}~: perte de contraste locale, détections manquées dans 15\% des cas
\end{itemize}

\hspace{0.65 cm}Ces limites sont intrinsèques à la tâche et partagées par tous les modèles comparés, comme documenté par Yang et al. \cite{Yang2016WIDERFACE}. Des stratégies de pré-traitement d'image (normalisation adaptative, rehaussement de contraste) ou d'augmentation de données ciblée pourraient atténuer ces faiblesses.

\subsubsection{Dépendance au dataset d'entraînement}

\hspace{0.65 cm}Les performances terrain (98\% Recall en classe) sont légèrement inférieures au mAP WIDER FACE (80\%). Cette différence s'explique par~:
\begin{itemize}
\item Des conditions d'éclairage spécifiques aux salles de classe (néons, fenêtres latérales) sous-représentées dans WIDER FACE
\item Des poses et expressions différentes (étudiants assis vs. sujets debout dans WIDER FACE)
\end{itemize}

\hspace{0.65 cm}Une phase de fine-tuning sur données de classe collectées in situ pourrait améliorer la généralisation, comme suggéré par Woo et al. \cite{Woo2018CBAM}.

\subsubsection{Latence sur hardware bas de gamme}

\hspace{0.65 cm}Les tests de latence ont été réalisés sur GPU NVIDIA A100 (benchmark) et smartphone haut de gamme (Snapdragon 888). Sur processeurs ARM plus anciens (Snapdragon 660, MediaTek Helio), la latence augmente à 15-20ms, limitant la fluidité temps réel. Une optimisation via quantification INT8 ou pruning structuré pourrait réduire cette latence, au prix d'une légère perte de précision (<1 point mAP).

\subsubsection{Équilibrage des coefficients de fusion attentionnelle}

\hspace{0.65 cm}L'intégration des sorties ECA et CBAM SAM utilise actuellement une fusion par produit élément par élément. Wang et al. \cite{Wang2024ECACBAM} suggèrent l'introduction de coefficients de pondération apprenables $\alpha$ et $\beta$ pour adapter dynamiquement l'influence de chaque mécanisme. L'exploration systématique de cette piste (grid search ou apprentissage par renforcement) pourrait apporter un gain additionnel de 0.5-1.0 point mAP.

\subsection{Perspectives d'amélioration}

\hspace{0.65 cm}Plusieurs axes d'amélioration émergent de cette étude~:

\subsubsection{Court terme}

\begin{itemize}
\item \textbf{Collecte et fine-tuning sur données de classe}~: acquisition de 2000-5000 images annotées en environnement scolaire réel pour améliorer la généralisation terrain
\item \textbf{Optimisation mobile avancée}~: quantification INT8, pruning, export ONNX/TensorRT pour latence <5ms sur hardware bas de gamme
\item \textbf{Coefficients adaptatifs}~: intégration de coefficients $\alpha$, $\beta$ apprenables pour fusion attentionnelle optimale
\end{itemize}

\subsubsection{Moyen terme}

\begin{itemize}
\item \textbf{Extension multi-benchmark}~: évaluation sur FDDB, AFW, MAFA (masques) pour valider la généralisation internationale
\item \textbf{Attention temporelle}~: intégration d'un module LSTM ou Transformer pour exploiter la cohérence temporelle en vidéo
\item \textbf{Détection multi-tâche}~: ajout de branches pour estimation d'âge, genre, expression faciale (réutilisant le backbone commun)
\end{itemize}

\subsubsection{Long terme}

\begin{itemize}
\item \textbf{Architectures de nouvelle génération}~: exploration de Vision Transformers légers (MobileViT, EfficientFormer) intégrant ECA-CBAM
\item \textbf{Apprentissage auto-supervisé}~: pré-entraînement sur millions d'images non annotées pour améliorer les représentations apprises
\item \textbf{Déploiement edge computing}~: portage sur microcontrôleurs (STM32, ESP32) pour applications IoT à très faible consommation
\end{itemize}

\subsection{Synthèse et contributions}

\hspace{0.65 cm}En synthèse, ce travail apporte trois contributions principales~:

\begin{enumerate}
\item \textbf{Architecture hybride ECA-CBAM}~: Première intégration documentée d'ECA-Net et CBAM SAM dans un détecteur facial mobile léger, atteignant l'état de l'art sur WIDER FACE (80.0\% mAP \textit{hard}) avec latence réduite (3.2ms)

\item \textbf{Méthodologie d'entraînement multi-phase}~: Stratégie d'activation séquentielle des mécanismes d'attention garantissant convergence stable et équilibre canal/spatial, généralisable à d'autres architectures hybrides

\item \textbf{Validation applicative}~: Démonstration de robustesse en contexte réel (98\% Recall en salle de classe) et analyse critique des limites, ouvrant des pistes concrètes pour recherches futures
\end{enumerate}

\hspace{0.65 cm}Ces contributions s'inscrivent dans la continuité des travaux de Kim et al. \cite{Kim2025FeatherFace}, Wang et al. \cite{Wang2020ECANet, Wang2024ECACBAM}, et Woo et al. \cite{Woo2018CBAM}, tout en proposant une synthèse originale répondant aux contraintes spécifiques de la détection faciale mobile pour vérification de présence en milieu éducatif.

\end{onehalfspace}

%\chapter{Implications sur le Système Éducatif}
\begin{onehalfspace}
\hspace{0.65cm}Les chapitres précédents ont étalé la conception et la réalisation de l’application \gls{scorq}-EDU destinée aux élèves et aux enseignants. Le présent chapitre montre les avantages qu’apporte son utilisation dans le processus d’enseignement/apprentissage tant chez les apprenants que chez les enseignants et sur le système éducatif Camerounais.
\section{Apport de \gls{scorq}-EDU pour les élèves }

\hspace{0.65cm}\gls{scorq}-EDU sera un outil d'apprentissage qui va susciter le gout d'apprendre chez l'apprenant. Les différents exercices que l’enseignant va publier permettront d'évaluer à chaque fois ses performances à travers soit les erreurs qu'il aura commises, soit alors d'évoluer en cas de réussite. De même, l'apprenant aura non seulement la motivation pour apprendre mais aussi la capacité de développer ses propres connaissances et de faire les révisions de ses leçons. L’opportunité de faire des raisonnements inductifs, d’augmenter les habiletés visuelles.
\section{Apport de SCORQ-EDU pour les enseignants}
\hspace{0.65cm}Pour l'enseignant SCORD-EDU lui permettra de déposer   tous les exercices que l'apprenant pourra s'exercer et de voir par la suite les listes des participants et même le suivi à distance de ces apprenants.
\section{Apport de SCORQ-EDU dans le système éducatif camerounais}
\hspace{0.65cm}SCOR-EDU met le joueur-apprenant au cœur de ses apprentissages, il favorise le développement des compétences comme la collaboration et la communication, la capacité à résoudre des problèmes, la créativité, la pensée critique et l’initiative. Cela appuie approche pédagogique dit APC et contribue à l’insertion des TIC en milieu éducatif camerounais.
\end{onehalfspace} 
\chapter*{Conclusion Générale et Perspectives}
\addcontentsline{toc}{chapter}{Conclusion Générale et Perspectives}
\begin{onehalfspace}

\hspace{0.65 cm}Au terme de ce mémoire, consacré à la conception, l'implémentation et l'évaluation d'une architecture hybride de détection faciale mobile optimisée par les mécanismes d'attention ECA-CBAM, il convient de dresser un bilan des contributions apportées, des résultats obtenus, et des perspectives ouvertes par ce travail.

\section*{Rappel du contexte et de la problématique}

\hspace{0.65 cm}La détection faciale sur dispositifs mobiles représente un défi technique majeur, conjuguant des contraintes antagonistes~: précision élevée sur petits visages et conditions difficiles d'une part, légèreté computationnelle et rapidité d'inférence d'autre part. Les applications pratiques, notamment la vérification de présence automatisée en salle de classe via smartphone, nécessitent des modèles capables de traiter en temps réel des images capturées à distance, avec éclairage variable, occlusions partielles, et groupes denses de personnes.

\hspace{0.65 cm}Les architectures légères existantes (BlazeFace, MobileFaceNets) privilégient la vitesse au détriment de la précision, tandis que les modèles lourds (RetinaFace) atteignent d'excellentes performances mais demeurent inadaptés au déploiement mobile. FeatherFace CBAM \cite{Kim2025FeatherFace} a récemment proposé un compromis prometteur (78.3\% mAP \textit{hard}, 0.49M paramètres), mais son module d'attention canal présente une complexité computationnelle de $O(C^2)$ limitant l'efficacité sur smartphone.

\hspace{0.65 cm}Cette problématique a guidé l'ensemble de notre démarche~: comment optimiser l'attention canal tout en conservant la robustesse spatiale, pour améliorer simultanément précision et efficacité sur mobile~?

\section*{Synthèse des objectifs atteints}

\hspace{0.65 cm}Ce travail s'est articulé autour de cinq objectifs spécifiques, tous atteints avec succès~:

\subsection*{1. État de l'art approfondi et critique}

\hspace{0.65 cm}Une revue exhaustive de la littérature a été conduite, couvrant~:
\begin{itemize}
\item Les architectures de détection faciale légères (Single-Shot Detectors), de Viola-Jones \cite{ViolaJones2001} aux modèles profonds récents
\item Les mécanismes d'attention en vision par ordinateur, avec analyse détaillée d'ECA-Net \cite{Wang2020ECANet} et CBAM \cite{Woo2018CBAM}
\item Les modèles mobiles spécialisés (BlazeFace, MobileFaceNets, FeatherFace) et leurs limitations respectives
\item Le benchmark WIDER FACE \cite{Yang2016WIDERFACE} et les protocoles d'évaluation standard
\end{itemize}

\hspace{0.65 cm}Cette revue a permis d'identifier précisément le point d'amélioration optimal~: le module d'attention canal de CBAM, dont le remplacement par ECA-Net constitue l'innovation centrale de ce mémoire.

\subsection*{2. Conception de l'architecture hybride FeatherFace ECA-CBAM}

\hspace{0.65 cm}L'architecture proposée intègre de manière originale~:
\begin{itemize}
\item Un backbone MobileNet-0.25 pour extraction de caractéristiques légères
\item Des modules hybrides ECA-CBAM à deux niveaux (backbone et BiFPN), combinant attention canal efficace ($O(C)$) et attention spatiale robuste
\item Un BiFPN pour fusion multi-échelle bidirectionnelle \cite{Tan2020EfficientDet}
\item Des têtes de détection multitâches (classification, régression de boîtes, landmarks)
\end{itemize}

\hspace{0.65 cm}Cette hybridation, jamais documentée auparavant dans le contexte de la détection faciale mobile, constitue une contribution architecturale validée empiriquement.

\subsection*{3. Développement d'une méthodologie d'entraînement multi-phase}

\hspace{0.65 cm}Une stratégie d'entraînement en trois phases séquentielles a été élaborée, inspirée de Wang et al. \cite{Wang2024ECACBAM}~:
\begin{enumerate}
\item Pré-entraînement du backbone sans attention (30 epochs)
\item Activation séquentielle ECA puis CBAM SAM (50 epochs)
\item Fine-tuning global avec régularisation intensive (30 epochs)
\end{enumerate}

\hspace{0.65 cm}Cette méthodologie garantit la convergence stable et évite la dominance d'un mécanisme attentionnel au détriment de l'autre, problème récurrent dans les architectures hybrides. Elle constitue une contribution méthodologique transférable à d'autres contextes d'apprentissage multi-module.

\subsection*{4. Évaluation rigoureuse sur benchmarks et contexte réel}

\hspace{0.65 cm}Le modèle a été évalué selon un protocole rigoureux~:
\begin{itemize}
\item Sur WIDER FACE \cite{Yang2016WIDERFACE}, benchmark de référence avec trois niveaux de difficulté
\item En contexte applicatif réel (salle de classe, 500 images terrain)
\item Via tests d'ablation systématiques (4 variantes architecturales)
\item En comparaison avec 6 modèles de l'état de l'art
\end{itemize}

\hspace{0.65 cm}Les métriques mesurées (mAP, Precision, Recall, latence, taille mémoire) suivent les standards de la communauté, garantissant la reproductibilité et la validité scientifique des résultats.

\subsection*{5. Analyse critique et positionnement scientifique}

\hspace{0.65 cm}Une discussion approfondie a été menée, incluant~:
\begin{itemize}
\item Interprétation des gains de performance (+1.7 points mAP \textit{hard}, -20\% latence)
\item Comparaison critique avec la littérature (FeatherFace CBAM, BlazeFace, MobileFaceNets, RetinaFace)
\item Identification de 4 limites majeures (cas extrêmes, dépendance dataset, hardware bas de gamme, coefficients fusion)
\item Proposition de 9 axes d'amélioration court/moyen/long terme
\end{itemize}

\hspace{0.65 cm}Cette analyse critique situe précisément la contribution dans le paysage scientifique et ouvre des perspectives concrètes pour recherches futures.

\section*{Principaux résultats et contributions}

\hspace{0.65 cm}Ce travail apporte trois contributions majeures, validées expérimentalement~:

\subsection*{Contribution 1~: Architecture hybride ECA-CBAM pour détection mobile}

\hspace{0.65 cm}Le modèle FeatherFace ECA-CBAM atteint des performances inédites pour un détecteur facial mobile ultra-léger~:

\begin{itemize}
\item \textbf{mAP WIDER FACE}~: 92.5\% (\textit{easy}), 90.8\% (\textit{medium}), \textbf{80.0\%} (\textit{hard})
\item \textbf{Latence d'inférence}~: 3.2 ms/image (GPU A100), 10 ms/image (smartphone Snapdragon 888)
\item \textbf{Taille mémoire}~: 1.4 MB (0.45M paramètres)
\item \textbf{Recall terrain}~: 98\% en contexte salle de classe avec <2\% faux positifs
\end{itemize}

\hspace{0.65 cm}Ces résultats surpassent tous les modèles mobiles légers comparés, tout en approchant à 98\% les performances de RetinaFace (modèle lourd de référence, >100 MB) avec un modèle \textbf{71× plus compact} et \textbf{19× plus rapide}.

\hspace{0.65 cm}L'hybridation ECA-CBAM démontre empiriquement que l'efficacité paramétrique d'ECA-Net \cite{Wang2020ECANet} (réduction $O(C^2) \rightarrow O(C)$) et la robustesse spatiale de CBAM \cite{Woo2018CBAM} se renforcent mutuellement, validant l'hypothèse centrale du mémoire.

\subsection*{Contribution 2~: Méthodologie d'entraînement multi-phase optimisée}

\hspace{0.65 cm}La stratégie d'activation séquentielle des mécanismes d'attention garantit~:
\begin{itemize}
\item Convergence rapide et stable (80\% mAP \textit{hard} en 100 epochs vs. >150 epochs pour entraînement joint naïf)
\item Équilibre optimal canal/spatial (ablation~: ECA seul 76.8\%, CBAM seul 78.3\%, hybride 80.0\%)
\item Absence de surapprentissage (courbes validation parallèles à entraînement)
\end{itemize}

\hspace{0.65 cm}Cette méthodologie, inspirée de Wang et al. \cite{Wang2024ECACBAM} mais adaptée au contexte de détection faciale, constitue un apport méthodologique généralisable à toute architecture combinant plusieurs modules d'attention complémentaires.

\subsection*{Contribution 3~: Validation applicative et analyse critique}

\hspace{0.65 cm}Au-delà des benchmarks académiques, ce travail~:
\begin{itemize}
\item Démontre la viabilité pratique pour vérification de présence en classe (98\% Recall, latence 10ms smartphone)
\item Identifie précisément les limites techniques (masques >70\%, visages <5px, contre-jour)
\item Propose des pistes d'amélioration concrètes (quantification INT8, coefficients adaptatifs, fine-tuning terrain)
\item Ouvre des perspectives à court, moyen et long terme (extension multi-benchmark, attention temporelle, Vision Transformers légers)
\end{itemize}

\hspace{0.65 cm}Cette démarche d'analyse critique rigoureuse distingue ce travail des simples études de performance, en apportant une réflexion approfondie sur les conditions de déploiement réel et les voies d'amélioration future.

\section*{Limites et perspectives}

\hspace{0.65 cm}Malgré les résultats prometteurs, plusieurs limites doivent être reconnues, ouvrant autant de perspectives de recherche~:

\subsection*{Limites identifiées}

\begin{enumerate}
\item \textbf{Sensibilité aux cas extrêmes}~: Masques couvrant >70\% du visage (68\% détection), visages ultra-petits <5px (45\% précision), contre-jour fort (15\% échecs). Ces limites sont intrinsèques à la tâche mais pourraient être atténuées par pré-traitement adaptatif ou augmentation de données ciblée.

\item \textbf{Dépendance au dataset d'entraînement}~: Performances légèrement inférieures en contexte réel (98\% Recall classe vs. 80\% mAP WIDER FACE) dues aux spécificités d'éclairage et de pose. Un fine-tuning sur 2000-5000 images terrain améliorerait la généralisation.

\item \textbf{Latence sur hardware ancien}~: Sur processeurs ARM anciens (Snapdragon 660), latence de 15-20ms limitant la fluidité. Quantification INT8 ou pruning structuré réduiraient cette latence au prix de <1 point mAP.

\item \textbf{Coefficients de fusion à optimiser}~: Fusion actuelle par produit élément par élément. Introduction de coefficients apprenables $\alpha$, $\beta$ pourrait apporter +0.5-1.0 point mAP.
\end{enumerate}

\subsection*{Perspectives à court terme (6-12 mois)}

\begin{itemize}
\item Collecte de données terrain annotées (2000-5000 images salle de classe) et fine-tuning spécifique
\item Optimisation mobile avancée~: quantification INT8, pruning, export ONNX/TensorRT/CoreML
\item Intégration de coefficients adaptatifs pour fusion attentionnelle optimale
\item Tests sur smartphones bas/milieu de gamme pour validation accessibilité
\end{itemize}

\subsection*{Perspectives à moyen terme (1-2 ans)}

\begin{itemize}
\item Extension multi-benchmark~: évaluation sur FDDB, AFW, MAFA (masques) pour généralisation internationale
\item Attention temporelle~: module LSTM ou Transformer léger pour cohérence vidéo, réduisant faux positifs
\item Détection multi-tâche~: branches additionnelles (âge, genre, expression) réutilisant backbone commun
\item Déploiement industriel~: intégration dans application mobile complète avec interface utilisateur
\end{itemize}

\subsection*{Perspectives à long terme (2-5 ans)}

\begin{itemize}
\item Architectures de nouvelle génération~: Vision Transformers légers (MobileViT, EfficientFormer) intégrant ECA-CBAM
\item Apprentissage auto-supervisé~: pré-entraînement sur millions d'images non annotées (SimCLR, DINO) pour représentations plus riches
\item Déploiement edge computing~: portage microcontrôleurs (STM32, ESP32, Raspberry Pi) pour IoT très basse consommation
\item Extension domaines connexes~: sécurité, santé (détection fatigue/émotion), retail (analyse client)
\end{itemize}

\section*{Apport scientifique et impact}

\hspace{0.65 cm}Ce mémoire s'inscrit pleinement dans la dynamique actuelle de recherche sur la détection faciale mobile, en proposant une synthèse originale des travaux récents de Kim et al. \cite{Kim2025FeatherFace}, Wang et al. \cite{Wang2020ECANet, Wang2024ECACBAM}, et Woo et al. \cite{Woo2018CBAM}.

\hspace{0.65 cm}L'apport scientifique se situe à trois niveaux~:

\begin{enumerate}
\item \textbf{Architecture}~: Première intégration documentée d'ECA-Net et CBAM SAM dans un détecteur facial mobile, avec validation empirique de la synergie canal/spatial

\item \textbf{Méthodologie}~: Stratégie d'entraînement multi-phase généralisable, garantissant convergence et équilibre dans architectures hybrides

\item \textbf{Application}~: Démonstration de viabilité pratique pour vérification de présence en classe, domaine sous-exploré dans la littérature académique
\end{enumerate}

\hspace{0.65 cm}L'impact potentiel est double~:

\begin{itemize}
\item \textbf{Scientifique}~: Ouverture d'une voie de recherche (attention hybride efficace) transférable à d'autres tâches de vision mobile (détection d'objets, segmentation, tracking)

\item \textbf{Sociétal}~: Facilitation de solutions technologiques accessibles (smartphone standard) pour améliorer la qualité des services éducatifs, sans infrastructure coûteuse
\end{itemize}

\section*{Mot de la fin}

\hspace{0.65 cm}En conclusion, ce travail valide l'intérêt de l'hybridation ECA-CBAM pour la détection faciale mobile, atteignant l'état de l'art sur WIDER FACE (80.0\% mAP \textit{hard}) avec efficacité computationnelle accrue (3.2ms latence, 1.4MB taille). La méthodologie d'entraînement multi-phase proposée garantit convergence et équilibre, tandis que la validation terrain (98\% Recall en classe) démontre la viabilité applicative.

\hspace{0.65 cm}Au-delà des résultats quantitatifs, ce mémoire illustre la pertinence d'une démarche scientifique rigoureuse~: revue critique de l'existant, hypothèse d'innovation clairement formulée, protocole expérimental reproductible, analyse critique des limites, et perspectives concrètes pour recherches futures.

\hspace{0.65 cm}Les perspectives ouvertes, à court, moyen et long terme, confirment le potentiel de cette approche pour contribuer à la nouvelle génération de modèles de vision embarquée, conjuguant précision, efficacité et accessibilité.

\hspace{0.65 cm}Ce travail s'achève donc sur une contribution solide à l'état de l'art de la détection faciale mobile, tout en ouvrant de nombreuses voies d'amélioration et d'extension pour les recherches à venir, tant du point de vue algorithmique que de celui du transfert vers des applications réelles à fort impact sociétal.

\end{onehalfspace}

%\chapter*{Bibliographie}
\addcontentsline{toc}{chapter}{Bibliographie}
\bibliographystyle{apalike}
\bibliography{Bibliographie/BibliographieBD} 
%\appendix
\newpage	\pagestyle{fancy}
	\renewcommand{\headrulewidth}{1pt}

	\fancyhead[R]{ANNEXES}
	\fancyhead[L]{}
	\fancyfoot[C]{\thepage}
	\renewcommand{\footrulewidth}{1pt}
	\fancyfoot[L]{\footnotesize{FOKOU Arnaud}}
	\fancyfoot[R]{\copyright Mémoire de Master II en Informatique}
%\newpage \pagenumbering{roman}%numérotation avec les chiffres romains
%\chapter*{Annexes}
\addcontentsline{toc}{chapter}{Annexes}
%\thispagestyle{empty}

\section*{Annexe 1 : Questionnaire adressé aux enseignants}
\addcontentsline{toc}{section}{Annexe 1 : Questionnaire adressé aux enseignants}

\subparagraph{Dans l'optique de conçevoir un outil informatique (LMS pour Learning Management System) qui vise à améliorer le système éducatif camerounais basé sur l'approche par compétences dans les établissements secondaires, nous venons auprès de vous pour collecter un ensemble d'information sur le fonctionnement de ce système en environnement réel. Votre contribution à la réalisation de ce travail nous est indispensable. A cet effet, nous vous prions de répondre aux questions suivantes en toute honnêteté et sans appréhension. Par ailleurs, nous garantissons votre anonymat.\\\\}

\textbf{I- Identification de l'enseignant}
\begin{enumerate}
  \item Quel est votre grade ?  PCEG.........	PLEG......... Contractuel........  Vacataire.........
  \item Dans quel établissement enseignez-vous ? ...................................................................
  \item Quel est votre nombre d'années de service ? ...........................................
  \item Quel est votre âge ? .........
  \item Dans combien d’établissements avez-vous déjà enseigné ? ...............................
  \item Quels sont les compétences visées par ces leçons ? .............................................
	\\
	
	\textbf{II- Enseignement des leçons portant sur \emph{«La programmation en HTML»}}\\
	
  \item Pendant combien d'années avez-vous déjà enseigné ces leçons ? ................................
  \item Quelles sont les compétences visées par ces leçons ?\\ ................................................................................................................... \\................................................................................................................... \\
................................................................................................................... \\................................................................................................................... \\
  \item Quelle(s) méthode(s) utilisez-vous pendant vos enseignements ?\\\\
Questionnement ..........  Magistrale/Expositive ........... Expérimentale .............\\
Autres : lesquelles ?.............................................................
  \item Vous enseignez ces leçons de manière :\\
	Théorique ..........  Pratique...........  Les deux .............
  \item Quelles sont les activités pratiques prévues liées à l'enseignement de ces leçons ?\\ ................................................................................................................... \\.....................................................................................................................\\

\textbf{III- Difficultés rencontrées pendant le processus d'enseignement-apprentissage des leçons liées au \emph{« La programmation en HTML »}.}\\

\item Les élèves ont-ils les mêmes difficultés dans différents établissements ?\\
 OUI ..........  NON ...........
\item Comment se comportent généralement les élèves pendant le passage de la phase théorique de ces leçons ?\\
Ils sont attentifs ..........			Ils ne sont attentifs ..........\\
Ils répondent aux questions ...........	Ils ne répondent pas aux questions ...........\\
Ils posent les questions ..........	Ils ne posent pas les questions ............
 \item Comment se comportent généralement les élèves pendant le passage de la phase pratique de ces leçons ?\\
Ils sont attentifs ..........			Ils ne sont attentifs ..........\\
Ils répondent aux questions ...........	Ils ne répondent pas aux questions ...........\\
Ils posent les questions ..........	Ils ne posent pas les questions ............
\item Vos élèves ont-ils des difficultés particulières d'assimilation/de compréhension de ces leçons ?\\
OUI..........  			NON.............\\
Si OUI, lesquelles ? ....................................................
\item De quel matériel avez-vous besoin pour convenablement conduire la(les) phase(s) pratique(s) de ces leçons ?\\
VIDEOPROJECTEUR.......      PLANCHES.......    ORDINATEURS.......  AUTRE....
\item Avez-vous à votre disposition tout le matériel nécessaire pour effectuer la pratique ?\\
OUI............ 	NON .............\\
Si NON, Quels matériels manquent ? ........................
\item Selon vous, quel le pourcentage des élèves capable de donner la structure générale d'un code HTML ?\\
Moins de 25\% ..... Entre 25\% et 50\% ....... 	Entre 50\% et 75\% ..... Plus de 75\% ......
\item Selon vous, quel pourcentage de vos élèves est capable de citer quelques balises et attributs en HTML ?\\
Moins de 25\% ..... Entre 25\% et 50\% ....... 	Entre 50\% et 75\% ..... Plus de 75\% ......\\

\textbf{IV.	Utilisation d'un outil TIC d'aide au processus d'enseignement-apprentissage des leçons portant sur \emph{« La programmation en HTML »}}\\

\item Pensez-vous qu'en utilisant un outil TIC pour l'apprentissage, vos élèves produiront de meilleurs résultats ?\\	OUI .......... NON ..............	
\item Avez-vous déjà utilisé un (des) outil(s) TIC d'aide à l'apprentissage pour enseigner ?\\
OUI .......... NON ..............
a) Si OUI, Lesquels ? ...............................................

\item Utilisez vous des cours disponible en ligne lors de la préparations de vos leçons?\\
    OUI.......... NON...........
\item Après l'utilisation de cet (ces) outil(s) TIC, avez-vous constatez une augmentation des performances de vos élèves ?\\ 	
OUI .......... NON .............. 		
\item Jugez-vous utile d'utiliser une plateforme éducative interactive pour l'enseignement des leçons portant \emph{«La programmation en HTML»} ?\\ 		
OUI .......... NON ..............\\
Si oui, quelles sont vos attentes par rapport à un tel outil ?\\
Qu'il intègre des simulations .......................\\
Qu’il intègre des médias audio et vidéo .............\\
Autres attentes : ...............................................................\\

\emph{\textbf{Nous vous remercions pour votre aimable contribution !!!}}
\end{enumerate}

\newpage
\section*{Annexe 2 : Questionnaire adressé aux élèves}
\addcontentsline{toc}{section}{Annexe 2 : Questionnaire adressé aux élèves}

\subparagraph{Dans le cadre de la fin de notre formation au Département d'Informatique et de Technologies Educatives de l'École Normale de l'Université de Yaoundé 1, nous menons une étude qui porte sur l'« analyse, conception et réalisation d'un module LMS (Learning Management System) prenant en compte L'APC utilisée dans les établissements secondaires». Votre contribution à la réalisation de ce travail nous est indispensable. A cet effet, nous vous prions de répondre aux questions suivantes en toute honnêteté et sans appréhension. Par ailleurs, nous garantissons votre anonymat.\\\\}

\textbf{I- Identification de l'élève}
\begin{enumerate}
\item Dans quel établissement fréquentes-tu ?.......................................
\item Es-tu nouveau ou ancien dans la classe de PA(ESP ou ALL) ?	Nouvel .....	Ancien .......
\item Quel âge as-tu ? .....................
\item Genre : Garçon ....... Fille .........\\

\textbf{II-	Apprentissage des leçons portant sur \emph{«La programmation en HTML»}}\\

\item Comment trouves-tu les leçons portant sur « La programmation en HTML » ?
Très difficiles....... Difficiles ...... Faciles ....... Très faciles .......
\item Par rapport aux leçons qui portent sur « La programmtion en HTML »,\\
a)	Les mots clés de ces leçons étaient faciles à comprendre\\
OUI ...... NON .........\\
b)	Les schémas étaient faciles à comprendre et à interpréter\\
OUI ...... NON .........\\
c)	L'enseignant posait des questions pendant les leçons\\
OUI ...... NON .........\\
d)	Je répondais aux questions posées par l'enseignant\\
OUI ...... NON .........\\
e)	Tu as fait des exercices d'application avec l'enseignant pendant les leçons\\
OUI ...... NON .........\\
\item Tu te sers de quel(s) outil(s) pour apprendre la programmation en HTML en plus du cours\\
Les planches ....... Internet ...... Le livre au programme .........\\ 	
Les bords ........	 Aucun de tous ......... Autres ........
\item Ton enseignant d'informatique utilise-t-il des outils TIC (didacticiel, vidéoprojecteur, animations, …) pour enseigner ?\\ 	
OUI........ NON ..........
\item As-tu l'habitude de travailler sur un ordinateur ou un Smartphone (Ex : Téléphone Android) ?\\
OUI ......... NON ..........
\item Quel est ton nombre d'heure d’apprentissage de tes leçons d'informatique par semaine ?\\
30 min ..... 1h ...... 1h30 min ....... 2h ....... plus de deux heures ..........
\item Quelle(s) activité(s) préfères-tu pendant tes temps libres ?\\
Manipuler le téléphone ......	Manipuler l'ordinateur ........ Lire ......\\
Visionner .......... Ecouter de la musique ..........
\item As-tu l'habitude de naviguer sur Internet ?\\
OUI......  NON .......

\textbf{III-	Utilisation d'un outil TIC d'aide au processus d'apprentissage des leçons portant sur « La programmation en HTML »}\\

\item Quel(s) outil(s) informatiques utilises-tu pour mieux comprendre les leçons portant sur « La programmation en HTML » ?\\	
Les Ordinateurs ....... SmartPhone ....... Tablette ...... Autre .......
\item As-tu déjà utilisé un logiciel éducatif (notamment un didacticiel, une plateforme éducative) pour apprendre ?\\
OUI ....... NON .......\\
Si OUI, cette utilisation t’a-t-elle permise de produire de meilleurs résultats en classe ?	
OUI ....... NON .......
\item Penses-tu qu'utiliser un didacticiel pour apprendre les  leçons portant sur « La programmation en HTML » peut augmenter tes performances ?\\
OUI ...... NON .......
\item Quelle est ta couleur préférée ?\\
Bleu ........ Jaune ......... Rouge .........  Vert ..........\\
Orange ....... Noir ......... Blanc..........  Autre(s).......
\item Notre intention est de développer une plateforme pour vous aider à mieux apprendre les leçons de votre classe. Qu'aimerez-vous le plus avoir dans cette plateforme ?\\
Notes de cours ...........\\
Exercices ................\\
Animations ...............\\
Vidéos....................\\
Jeux portant sur les leçons...........

\item Comment apprenez- vous le mieux ?\\
Par des explications ............\\
En pratiquant ...................\\
En regardant des images ............\\
En lisant ...............\\

	\emph{\textbf{Nous te remercions pour ton aimable contribution !!!}}
\end{enumerate}

\newpage
\section*{Annexe 3 : Manuel d'utilisation du module LMS}
\addcontentsline{toc}{section}{Annexe 3 : Manuel d'utilisation de l'application }

\subsection*{Présentation du module}
\subsection*{Obtention et Installation du module}
\subsection*{Démarrage du module}
\subsection*{Prise en main du module} 
%-----------------Fin Zone 5: fin de la numérotation avec les chiffres arabs---------------------

\end{singlespace}
\end {document} 